% Modernised reading — fonds Alexandre Grothendieck, shelfmark 74
% (pages 1-135, the whole folder).
%
% This is an INTERPRETATION, not a transcription. It re-reads the pages in
% today's notation and vocabulary, states the hypotheses the manuscript leaves
% implicit, and drops the critical apparatus so that the mathematics reads
% straight through. Everything the brackets and underlines carried moves into
% a footnote. It is held to being mathematically correct as it stands; where
% the manuscript is loose or mistaken, the true statement is what appears here
% and the footnote says what the page has. In any dispute of substance the
% transcription governs: whoever cites Grothendieck must cite batch-01.fr.tex
% (pp. 1-20), batch-02.fr.tex (pp. 21-40), batch-03.fr.tex (pp. 41-60),
% batch-04.fr.tex (pp. 61-80), batch-05.fr.tex (pp. 81-100), batch-06.fr.tex
% (pp. 101-120) or batch-07.fr.tex (pp. 121-135).
%
% Pass: Opus 5.5 (claude-opus-5-5), 2026-10-10 — first pass, unchecked against the pages by a human.
% The folder was read whole, its seven batches together; this is one document
% for the shelfmark. It was made on Opus 5.5 and has not been measured against
% the reference reading of folder 115 (made on Opus 5). The literature named
% in the footnotes is cited from memory and was not checked against the
% sources. Model followed: the reading of folder 69 (same objects, same
% dating), cross-referenced throughout; folder 75 for torsors and omega(I).
% Checked by machine (Python, on the standard model of the 27 lines a_i, b_i,
% c_ij, on the duads of a 6-set, and on the 3x3 grid): 45 triangles, 720
% triads, 72 sixes, 120 pavés and the intersection pattern of p. 62 (1, 2,
% 27, 54, 36); the distribution 6/9/6/3 of p. 64; the models of pp. 51-52
% are srg(27,10,1,5); T_1(B) of p. 59 (6 independent 5-sets, each point in
% 2); the four orbits (18, 162, 216, 324) of the group of order 1296 on the
% triads (p. 60); |G°| = 36, 720, 25920 for c = 1, 2, 4 and G° transitive on
% ordered edges for c = 2, 4 (p. 107); the 45 double transpositions of S_6
% fixing exactly a line (p. 106); phi : Y -> E bijective for c = 1, 2-to-1 onto
% 4 points for c = 2, injective for c = 4 (pp. 127-129); the sign of the
% commutator formula of p. 33 (Heisenberg model); the arithmetic of pp.
% 115-120.
%
% Scope. Every transcribed page is used. Absent from the transcription as
% blank or without mathematics: pp. 2, 4, 7, 9, 11, 17, 19, 24, 31, 34, 36,
% 40, 42, 45 (covers with a pencilled number), 67 (typed Bourbaki seminar
% programme), 71 (blank separator), 91 (cover « (33 p) »), 126 (cover
% « (9p) »); p. 122 is a printed USTL examination cover, of which only his
% arithmetic counts; the upper half of p. 41 (household accounts) is not
% described. The leaves are partly out of order (p. 5 continues p. 6).
%
% Spine. A family of graphs — vertices, edges, triangles; every edge in one
% triangle; every vertex off a triangle linked to exactly one of its
% vertices — of which the graph of the 27 lines is the largest member:
%   pp. 1-16    ternary relations with unique completion (Latin squares),
%               bitorsors; the homomorphism lemma (pp. 5-6); the loop of a
%               pointed relation (pp. 8-12); associativity, a Moufang-type
%               identity (pp. 13-15); loops of order <= 4 (p. 16).
%   pp. 18-30   a graph built on a group (p. 18); the 27 in Schläfli's
%               notation xi_i, xi'_i, xi_ij: triads, hexagons, the counts 27,
%               36, 40, 45, 72, 120, 720; E_6 (p. 27); p. 28's Prop.
%   pp. 32-39   central extensions of Gamma = Ker(V^I -> V): z_rho, quadratic
%               maps, the universal module, the canonical extension for
%               dim V = 2 over F_2.
%   pp. 41-44   scratch on pentagons and the dodecahedron (720 = 6!).
%   pp. 46-60   the cubic graph: discrete and bi-discrete parts, triads
%               (linked, W-linked, bilinked), hexagons of triads, bitriangles,
%               the decomposition into three bitriangles (pp. 51-52), the
%               15-vertex graph T_1(B) (p. 59), the four types of triads.
%   pp. 61-70   bitriangles renamed « pavés »; the 120 pavés relative to one;
%               triad <-> hexagon; nu_tau; the pavé and its complement.
%   pp. 72-80   the diagram of coverings Tria' -> Hex -> triHex; hexagons as
%               (tau, eps, eps', phi); polygons n = 3, 4, 6 and the subgroup
%               lattices of D_4, D_6.
%   pp. 81-90   squares and hexagons in a pavé; products, coproducts and the
%               opposite graph; the pavé is self-opposite.
%   pp. 92-105  « complexes cubiques généralisés »: the axioms, the structure
%               around a triangle t_0, the Scholie, the reduction to (I) a
%               Latin square of order c and (II) double covers; counts in c.
%   pp. 106-116 automorphisms: the symmetries sigma_t, G°, the transitivity
%               theorem, lemmas, |G| for c = 1, 2; p. 116: elementary abelian.
%   pp. 117-123 the tower T_{-1} < T_0 < T_1 < T_2 (c = 0, 1, 2, 4): counts
%               of embedded and pinned substructures; categories of pairs.
%   pp. 121, 124-125 relations between the sigma_t; « short-cut ».
%   pp. 127-129, 135 a marked point: V = F_2^I, the map phi : Y -> E.
%   pp. 130-131 the same axioms in terms of triangles; configurations, duality.
%   pp. 132-135 « diagrams A_n », prisms, bicarrés, squares centred at a point.
%
% Notation fixed for the document.
%  - Delta = (S, A, T): vertices, edges, triangles; « liés » kept (for the 27
%    lines: the lines meet). Folder 69 writes C for the cubic graph; here
%    Delta, as on most of these pages.
%  - triangle = three pairwise linked vertices; triade = three pairwise
%    non-linked vertices; base (hexaplet) = six pairwise non-linked; bibase =
%    double-six. Pavé = bitriangle = the 3x3 grid (product of two triangles).
%  - c = card E_a (the number of triangles through a vertex, minus one);
%    T_nu for the graph with c = 2^nu (nu = 0, 1, 2) and T_{-1} the triangle
%    (c = 0). The pages write T, tau, T (cursive), always with this index.
%  - The 27: xi_i, xi'_i (i = 1..6), xi_ij, Schläfli's a_i, b_i, c_ij; eta the
%    class of a line of the plane, xi = sum xi_i; intersection form signs as
%    in folder 69 (xi_i.xi_i = -1).
%  - omega(I): the two circular orders on a 3-set (as in folders 69, 75).
%  Homonymies, stated in the spine section: « bitriade » (p. 28: 36 objects,
%  double-sixes; p. 46: a hexagon); « triades liées » (p. 46: union is a
%  base; p. 68: « parallèles » in a pavé); Gamma (a ternary relation, pp.
%  1-16; Ker(V^I -> V), pp. 32-39; c^2 points, pp. 101-103); I (a 3-set; the
%  basis of p. 38; the triangles at a, p. 127).
%
% Substantive departures, each footnoted where it occurs:
%  p. 10   sigma_2 sigma_1 = id etc. do not follow from a)-c); they are the
%          hypothesis u_1 u_2 u_3 = id admitted p. 8.
%  p. 14   margin « sigma(xy) = sigma x sigma y »: anti-automorphism (p. 13).
%  p. 15   the boxed identity at z = 1 is the left Moufang identity (ours).
%  p. 16   loops of order <= 4 are groups; the first non-associative loop has
%          order 5 (ours).
%  p. 18   « i in V_lambda » read I_lambda (transcription's doubt).
%  p. 26   (2 eta - xi).xi_j = 1 with the intersection form (page: 2).
%  p. 27   r_456 = r_1+2r_2+3r_3+2r_4+r_5+r_6 (page: 3(r_1+r_2+r_3)+...);
%          « indice 3 » holds for the lattice, not in W (index 240).
%  p. 33   z(x+y) - z(x) - z(y) = [psi_i(x), psi_{rho i}(y)] (page: the
%          opposite sign).
%  p. 35   over Z/n an extra condition (lambda(n e_i) = 0) is needed.
%  p. 38   « V + Gamma^2(V) » is the universal module of degree-<=2 maps
%          (Passi's I/I^3), an extension, not a direct sum; the result for
%          V = F_2^n is right.
%  p. 52   |Aut| = 6^4 = 1296 (page: « 3.6^4 », uncertain).
%  p. 72   Hex has 120 elements (page: 720).
%  p. 87   with p. 86's definition the coproduct of two segments is K_4; the
%          square drawn is their product.
%  p. 92   T ⊂ P_3(S) (page: P_3(A)).
%  p. 106  b) is false as stated (transpose of the grid; double transpositions
%          for c = 2); corrected characterisation.
%  p. 107  G° = G for c = 2 only; index 2 for c = 1 and for c = 4 (ours, by
%          machine for c = 4).
%  p. 119  c_{0,gamma} without the factor (c-1); p. 120's closed formulas
%          lack a factor c (both).
%  p. 123  the (0,2) entry: (Ens)_3 x (Ens)_{3,3} (order 432, as p. 52); the
%          second table's entries are doubtful and three disagree with the
%          counts.
%  p. 128  u is injective iff card I is even, i.e. nu = 0 (page: nu ≠ 0).
%
% Links the pages do not write (ours, said so in the body): the axioms of p.
% 92 are those of a generalised quadrangle of order (2, c), and c in {1, 2,
% 4} is the classical list; axiom I of p. 1 and beta_Gamma of p. 102 are the
% same Latin-square condition, so pp. 1-16 are problem (I) of p. 102; the
% extensions of pp. 32-39 and the diagram of p. 121 are problem (II); p. 116
% forces the F_2-vector space that p. 39 then extends; the 15-vertex T_1 is
% the Cremona-Richmond configuration; pavé = Steiner trihedral pair; the
% triangle stabiliser « Trialité » (p. 123) has order 1152 = |W(F_4)|.
% Folder 69 has the same family seen from the 27 (its pp. 2-16, 100-118).
%
% Announced and never established: p. 1's category C and the question « le
% foncteur phi est-il une équivalence ? » (p. 5); case c) of p. 18; p. 28's
% Prop (fragmentary); the extension of p. 37 « déterminé à isom. unique
% près » (stated); p. 63's programme a), b); the Prop of p. 86 (stated with
% « ? »); the Résolution of p. 100 (stops); the classification I), II) of
% pp. 102-103; the « corollaire plus fort » of p. 114 (struck); p. 129 b)
% (stops); the bracket of p. 135 (not closed).
\documentclass[11pt,a4paper]{article}
\input{../preamble/grothendieck}

\folder{74}
\pages{1}{135}
\dating{[à partir de 1976]}
\watermark{Édition de démonstration\\interprétation personnelle de l'œuvre}
\foldertitle{Complexes cubiques : notes manuscrites (s.d.) — lecture modernisée du dossier entier}

\begin{document}

\section*{Résumé}

\begin{resume}
Prenez trois points deux à deux reliés : c'est un triangle. Imaginez
maintenant un dessin fait de points et de traits où chaque trait appartient
à un seul triangle, et où tout point placé hors d'un triangle est relié à un
sommet de ce triangle et à un seul. La règle paraît anodine ; elle est en
réalité si contraignante qu'il n'existe qu'une poignée de tels dessins finis
non triviaux. Le plus petit a neuf points rangés en carré de trois sur trois,
deux points étant reliés quand ils sont sur une même ligne ou une même
colonne : Grothendieck l'appelle d'abord « bitriangle », puis « pavé ». Le
suivant a quinze points. Le plus grand en a vingt-sept, et c'est un objet du
XIX\textsuperscript{e} siècle : les vingt-sept droites d'une surface cubique,
deux droites étant reliées quand elles se coupent. Le dossier, intitulé
« Complexes cubiques », tourne tout entier autour de cette famille. Les
pages ne parlent jamais de la surface elle-même ; c'est le dessin des droites,
qu'elles appellent « graphe cubique », qui les occupe.

Le chemin n'est pas droit. Les premières pages sont de l'algèbre pure : trois
ensembles et une relation entre eux telle que deux éléments pris dans deux des
ensembles en déterminent un unique troisième --- autrement dit un carré latin,
une table où chaque valeur apparaît une fois par ligne et par colonne. Quand
une telle table est-elle celle d'un groupe ? Il montre qu'entre tables de
groupes les seules correspondances sont les homomorphismes, puis cherche ce
qui force l'associativité, et tombe sur une identité qu'on reconnaît
aujourd'hui comme une identité de Moufang. Viennent ensuite des pages de
dénombrements sur les vingt-sept droites --- 45 triangles, 72 « bases » de six
droites disjointes, 36 doubles-six, 120 pavés, 720 triplets de droites
disjointes --- avec le groupe de symétries d'ordre $51\,840$ et le système de
racines $E_6$ en arrière-plan, puis une étude minutieuse des configurations
qu'on trouve dans le graphe : triades, hexagones, pavés, et la façon dont
chacune en détermine d'autres.

Le tournant est au milieu du dossier. Il pose la règle du premier paragraphe
comme axiome, sans supposer qu'il y ait vingt-sept points, et regarde ce que
devient le graphe autour d'un triangle. Tout se ramène à deux questions, et
ce sont exactement les deux sujets des premières pages : un carré latin
(c'est l'algèbre du début) et un revêtement double de ce carré, qui relève de
la théorie des extensions de groupes, présente elle aussi dans le dossier. On
voit là, comme souvent chez lui, des pages d'apparence décousue qui se
laissent lire après coup comme les deux moitiés d'un même problème. Il calcule
alors, pour chaque taille possible, le nombre de triangles, de pavés, et
l'ordre du groupe des symétries, qu'il engendre par des « symétries » attachées
aux triangles ; les tableaux de la fin donnent $6$, $72$, $720$, $51\,840$.

Le dossier ne referme pas la classification qu'il annonce, mais tout y est
disposé pour elle. Les noms modernes à chercher sont ceux des quadrangles
généralisés d'ordre $(2, t)$ --- c'est le nom actuel de ces graphes, et l'on
sait qu'il n'y en a que pour $t = 1, 2, 4$ ---, des vingt-sept droites et des
doubles-six de Schläfli, des paires de trièdres de Steiner, du groupe de Weyl
de $E_6$, des carrés latins, quasigroupes et boucles de Moufang, des
extensions centrales et des formes quadratiques sur $\mathbb{F}_2$.

\keywords{generalized quadrangle, collinearity graph, 27 lines on a cubic surface, double-six, Steiner trihedral pair, Paley graph of order 9, Cremona-Richmond configuration, Weyl group W(E6), E6 root lattice, Latin square, 3-net, quasigroup, isotopy, Moufang loop, bitorsor, central extension, quadratic map, polynomial map of degree 2, Passi module, generated involutions, dihedral group subgroup lattice, complement of a graph, self-complementary graph}
\end{resume}

\pagerange{1}{135}
\section*{Le fil du dossier, et les conventions}

\subsection*{Les feuillets}

Le dossier compte cent trente-cinq pages au fac-similé. Plusieurs sont des
chemises ou des intercalaires qui ne portent qu'un nombre au crayon (pages 42,
45, 91 « (33 p) », 126 « (9p) ») ; d'autres sont blanches ou sans
mathématiques (2, 4, 7, 9, 11, 17, 19, 24, 31, 34, 36, 40, 71) ; la page 67
est un programme dactylographié du séminaire Bourbaki, la page 122 une copie
d'examen imprimée de l'Université des Sciences et Techniques du Languedoc, dont
seul compte un calcul écrit en travers ; le haut de la page 41 porte des
comptes domestiques. On ne décrit rien de tout cela davantage. Les feuillets
du début sont en partie dans le désordre : la page 5 continue une phrase
ouverte au bas de la page 6. La datation « à partir de 1976 » est celle des
archivistes ; elle s'accorde avec le dossier 69, qui traite des mêmes objets
et porte la même.

Le dossier ne contient presque aucun texte rédigé. C'est une suite de notes de
travail, souvent d'écriture très rapide, où les formules se lisent mieux que
la prose qui les relie ; on le signale chaque fois que la lecture qu'on
propose dépend d'un mot incertain.

\subsection*{Les stations}

\begin{itemize}
\item \textbf{Pages 1 à 16 : relations ternaires et carrés latins.} Trois
  ensembles et une partie $\Gamma$ de leur produit, deux coordonnées
  déterminant la troisième ; les bitorseurs ; le lemme d'homomorphisme
  (pages 5 et 6) ; la loi de composition d'une relation pointée (pages 8 à
  12) ; l'associativité et une identité de Moufang (pages 13 à 15) ; les
  petits ensembles (page 16).
\item \textbf{Pages 18 à 30 : premiers dénombrements.} Un graphe construit
  sur un groupe (page 18) ; les vingt-sept droites dans la notation de
  Schläfli, triades et hexagones, les nombres $27$, $36$, $40$, $45$, $72$,
  $120$, $720$ ; le diagramme $E_6$ (page 27).
\item \textbf{Pages 32 à 39 : extensions centrales.} Les relèvements
  $\psi_i$, la fonction $z_\rho$, les applications quadratiques, le module
  universel, l'extension canonique en dimension 2 sur $\mathbb{F}_2$.
\item \textbf{Pages 41 à 44 : un brouillon} sur les pentagones et le
  dodécaèdre.
\item \textbf{Pages 46 à 60 : le graphe cubique.} Parties discrètes et
  bi-discrètes, triades liées, hexagones de triades, bitriangles, la
  décomposition en trois bitriangles et ses modèles (pages 51 et 52), le graphe
  à quinze sommets (page 59), les quatre types de triades (page 60).
\item \textbf{Pages 61 à 70 : les pavés.} Les 120 pavés vus d'un pavé
  donné, la bijection triade--hexagone, la fonction $\nu_\tau$, le pavé et son
  complémentaire.
\item \textbf{Pages 72 à 80 : le diagramme des revêtements et les polygones.}
\item \textbf{Pages 81 à 90 : carrés et hexagones dans un pavé ; produits et
  coproduits de graphes.}
\item \textbf{Pages 92 à 105 : les complexes cubiques généralisés.} Les
  axiomes, la structure autour d'un triangle, la Scholie, la réduction de la
  classification à deux problèmes, les nombres en fonction de $c$.
\item \textbf{Pages 106 à 116 : automorphismes.} Les symétries
  $\sigma_{\underline t}$, le groupe $\mathcal{G}^\circ$ qu'elles engendrent,
  le théorème de transitivité.
\item \textbf{Pages 117 à 123 : la tour $\mathcal{T}_{-1} \subset
  \mathcal{T}_0 \subset \mathcal{T}_1 \subset \mathcal{T}_2$} et ses tableaux.
\item \textbf{Pages 121, 124, 125 : relations entre symétries.}
\item \textbf{Pages 127 à 129 et 135 : un point marqué}, l'espace
  $\mathbb{F}_2^I$.
\item \textbf{Pages 130 et 131 : les mêmes axiomes en termes de triangles.}
\item \textbf{Pages 132 à 135 : diagrammes $A_n$, prismes, carrés.}
\end{itemize}

Le dossier a donc trois couches : une algèbre (pages 1 à 16 et 32 à 39), une
géométrie combinatoire du graphe des vingt-sept droites (pages 18 à 90), et, à
partir de la page 92, une théorie axiomatique qui contient les deux premières.
Les pages 101 à 103 le disent presque en toutes lettres : classer les graphes
de la famille revient à classer (I) certains carrés latins et (II) leurs
revêtements doubles. Que le problème (I) soit celui des pages 1 à 16 et le
problème (II) celui des pages 32 à 39, les pages ne l'écrivent pas ; le
diagramme de la page 121, qui remet côte à côte $V$, $\Gamma$,
$\widetilde\Gamma$, $\lbrace \pm 1\rbrace$ et les $\widetilde E_i$, est ce qui
s'en approche le plus.

\subsection*{Conventions, valables pour tout le dossier}

\begin{itemize}
\item Un graphe est noté $\Delta = (S, A)$, ou $(S, A, T)$ quand on veut
  nommer l'ensemble $T$ de ses triangles ; $S$ est l'ensemble des sommets,
  $A \subset \mathfrak{P}_2(S)$ celui des arêtes. On garde le mot « liés » des
  pages. Le dossier 69, qui étudie le même graphe, le note $C$ ; on suit ici
  les pages de ce dossier-ci, qui l'écrivent $\Delta$ ou $S$.
\item Un \emph{triangle} est une partie de trois sommets deux à deux liés ;
  une \emph{triade}, une partie de trois sommets deux à deux non liés ; une
  \emph{base} (ou « hexaplet »), six sommets deux à deux non liés ; une
  \emph{bibase}, une paire de bases associées. Le \emph{pavé} (d'abord appelé
  « bitriangle ») est le graphe à neuf sommets $I \times J$, $\operatorname{card} I
  = \operatorname{card} J = 3$, où deux sommets distincts sont liés quand ils ont
  une coordonnée commune ; ses six triangles sont les trois lignes et les trois
  colonnes. Un \emph{hexagone} (« petit hexagone ») est un cycle induit de
  longueur 6.
\item Dans le graphe des vingt-sept droites, « liés » veut dire que les droites
  se coupent ; les triangles sont les 45 plans tritangents, les triades les
  triplets de droites deux à deux disjointes, les bases et bibases les « six »
  et les doubles-six de Schläfli, les pavés les paires de trièdres de Steiner.
  Ces noms classiques sont de nous ; le dossier ne nomme jamais la surface
  cubique, et l'identification repose sur les nombres qu'il calcule et sur le
  diagramme $E_6$ de la page 27\footnote{On a vérifié tous les nombres cités
  dans ce document sur le modèle standard (les droites $a_i$, $b_i$, $c_{ij}$,
  $a_i$ coupant $b_j$ pour $i \neq j$, $c_{ij}$ coupant $a_i, a_j, b_i, b_j$ et
  les $c_{kl}$ avec $\lbrace k, l\rbrace \cap \lbrace i, j\rbrace = \emptyset$).
  C'est aussi le modèle du dossier 69, vérifié de la même façon.}.
\item Les vingt-sept sont écrits, aux pages 20 à 30 et 124, $\xi_i$,
  $\xi'_i$ (ou $\hat\xi_i$, $\xi_i^*$, $\xi_i^\wedge$) et $\xi_{ij}$ ($1 \leq
  i, j \leq 6$), qui sont les $a_i$, $b_i$, $c_{ij}$ de Schläfli ; $\eta$ est
  la classe d'une droite du plan\footnote{Ces pages ne définissent ni $\eta$
  ni la forme ; on les lit avec les conventions du dossier 69 (pages 41 à
  44), où la surface est réalisée comme le plan éclaté en six points.}, $\xi = \sum \xi_i$, et la forme est la forme
  d'intersection, $\xi_i \cdot \xi_i = -1$, comme au dossier 69.
\item Dans la théorie générale (pages 92 et suivantes), $c$ désigne le cardinal
  de $\widetilde E_a / \sigma_a$, c'est-à-dire le nombre de triangles passant
  par un sommet, moins un. Les quatre graphes de la famille sont notés
  $\mathcal{T}_\nu$ avec $c = 2^\nu$ : $\mathcal{T}_0$ le pavé ($c = 1$, 9
  sommets), $\mathcal{T}_1$ le graphe à 15 sommets ($c = 2$), $\mathcal{T}_2$
  le graphe cubique ($c = 4$, 27 sommets), et $\mathcal{T}_{-1}$ le triangle
  seul ($c = 0$)\footnote{Les pages écrivent $\mathcal{T}_\nu$ (une capitale
  cursive de lecture douteuse, pages 57, 58, 63, 70, 75), $\tau_\beta$ (pages
  117 à 123) et $T_\nu$ (page 127), toujours avec cet indice. La page 123
  ajoute $\tau_{-\infty}$, le graphe vide, et écrit $c = 0 = 2^{-\infty}$ en
  regard du triangle, indexé pourtant $-1$ : l'indice et l'exposant ne
  coïncident qu'à partir de $\nu = 0$.}.
\item $\omega(I)$ est l'ensemble à deux éléments des ordres circulaires d'un
  ensemble $I$ à trois éléments, torseur sous $\lbrace \pm 1\rbrace$, comme aux
  dossiers 69 et 75 ; $X \wedge \omega$ le produit contracté.
\end{itemize}

Plusieurs homonymies traversent le dossier, et il faut les avoir en tête. Le
mot « bitriade » désigne page 28 des objets au nombre de 36, les bibases, et
page 46 un hexagone (deux triades dont la relation de non-liaison est une
bijection). « Deux triades liées » veut dire page 46 que leur réunion est une
base, et page 68 qu'elles sont « parallèles » dans un même pavé ; la page 68
le marque elle-même en renommant le bitriangle « pavé ». La lettre $\Gamma$
désigne une relation ternaire (pages 1 à 16), le groupe
$\operatorname{Ker}(V^I \to V)$ (pages 32 à 39) et l'ensemble à $c^2$
éléments des pages 101 à 103 --- trois objets qui, on le verra, sont trois
aspects d'une même chose. Enfin $I$ est un ensemble à trois éléments, sauf
page 38 (une base d'espace vectoriel) et page 127 (l'ensemble des triangles
passant par un sommet).

\pagerange{1}{3}
\section*{1. Relations ternaires et bitorseurs (pages 1 et 3)}

Soit $I$ un ensemble à trois éléments, $(E_i)_{i \in I}$ une famille
d'ensembles et $\Gamma \subset \prod_{i\in I} E_i$.

\textbf{Axiome I.} Si $I = \lbrace i, j, k\rbrace$, pour tous $x_i \in E_i$,
$x_j \in E_j$ il existe un unique $x_k \in E_k$ tel que $(x_i, x_j, x_k)
\in \Gamma$.

Autrement dit, les trois projections $\Gamma \to E_i \times E_j$ sont
bijectives : $\Gamma$ est un \emph{carré latin} (la table d'un quasigroupe,
ou, en géométrie, un \emph{3-réseau}), écrit symétriquement en ses trois
directions\footnote{Noms de nous. C'est la même condition que la condition (1)
du dossier 69 (page 107), et --- on le verra --- que l'axiome $\beta_\Gamma$
de la page 102 de ce dossier.}. Ces objets forment une catégorie
$\mathcal{C}$, $I$ variant, pour les isomorphismes au-dessus des bijections de
$I$\footnote{« ($I$ variable ; pour iso sur $I$) » : les deux derniers mots
sont incertains.}.

La seconde définition (\textbf{II}) donne la source des exemples. Soit
$(G_i)_{i\in I}$ une famille de groupes, et pour chaque $i$ un ensemble $E_i$
qui soit un \emph{bitorseur} sous les deux autres groupes : si $I = \lbrace
i, j, k\rbrace$, $G_j$ opère à gauche et $G_k$ à droite sur $E_i$, chacun
simplement transitivement, les deux actions commutant\footnote{La page
précise qu'on peut « renverser » l'une des actions en passant à l'inverse ;
une première ligne, biffée, parlait d'un système transitif d'isomorphismes
extérieurs entre les $G_i$, idée qui revient plus bas.}. Le produit $\prod_i
G_i$ opère alors diagonalement sur $P = \prod_i E_i$ par
\[
(g_i, g_j, g_k) \cdot (x_i, x_j, x_k) = (g_j x_i g_k^{-1},\ g_k x_j
g_i^{-1},\ g_i x_k g_j^{-1}).
\]
Pour $x \in P$, les trois conditions
\[
(\mathrm{a}_i) \quad \delta(G_i) x \subset \delta(G_j \times G_k) x
\]
(et leurs permutées circulaires) sont équivalentes\footnote{La ligne qui les
introduit est surchargée ; seul « équivalentes » se lit sûrement.}. Un point
$x$ définit, par $g_j x_i = x_i \tau(g_j)$, des isomorphismes
$\tau_{k,j}(x_i) : G_j \to G_k$, et la condition s'écrit
\[
\tau_{j,i}(x_k)\, \tau_{i,k}(x_j)\, \tau_{k,j}(x_i) = \mathrm{id}_{G_j},
\]
condition invariante par permutation circulaire. Sous cette condition l'orbite
$\Gamma$ de $x$ sous $\prod G_i$ est un torseur sous chacun des trois produits
$\prod_{j \neq i} G_j$, et elle vérifie l'axiome I : connaissant $a \in
\Gamma$ et $x_i$, $x_j$, on détermine successivement les éléments du groupe qui
mènent de $a$ à $x$, puis $x_k$. L'\emph{existence} d'un tel $x$ équivaut à ce
que les isomorphismes extérieurs (modulo les automorphismes intérieurs)
$\tilde\tau_{ij}$ définis par les bitorseurs forment un système transitif. On
obtient ainsi une catégorie $\mathcal{C}'$ et un foncteur $\varphi :
\mathcal{C}' \to \mathcal{C}$\footnote{Les lignes qui entourent l'encadré de
la page 3 sont d'une écriture très rapide et seuls les symboles s'y lisent ;
on rend le sens que les formules imposent. C'est la théorie des bitorseurs du
dossier 69 (page 112) ; elle relève de la cohomologie non abélienne de Giraud
(1971), que la page n'invoque pas.}.

\pagerange{5}{6}
\section*{2. Le lemme d'homomorphisme (pages 6 et 5)}

\textbf{Proposition} (page 6). Le foncteur $\varphi : \mathcal{C}' \to
\mathcal{C}$ est pleinement fidèle : il identifie $\mathcal{C}'$ à une
sous-catégorie pleine de $\mathcal{C}$.

On se ramène, les deux catégories étant fibrées sur le groupoïde des
ensembles à trois éléments, aux fibres en $I = \lbrace 1, 2, 3\rbrace$. La
fidélité est claire. Pour la plénitude, un point de $\Gamma$ identifie les
$G_i$ entre eux et aux $E_i$, donc l'objet au triple $(G, G, G)$ avec
\[
\Gamma = \lbrace (x, y, z) \in G^3 \mid xyz = 1\rbrace ;
\]
de même pour un second objet, avec un groupe $G'$, en prenant pour point
base l'image de $(1, 1, 1)$. Un morphisme de $\mathcal{C}$ est alors un
triple d'applications $u_1, u_2, u_3 : G \to G'$ avec $u_i(1) = 1$ et
\[
(*) \qquad xyz = 1 \ \Longrightarrow\ u_1(x)\, u_2(y)\, u_3(z) = 1 .
\]

\textbf{Lemme} (page 5). Sous ces conditions, $u_1 = u_2 = u_3$, et c'est un
homomorphisme de groupes.

En effet, $x = 1$ donne $u_2(y) u_3(y^{-1}) = 1$, et $y = 1$ donne $u_1(x)
u_3(x^{-1}) = 1$ ; donc $u_1(y) = u_2(y) = u_3(y^{-1})^{-1}$ pour tout $y$.
Avec $z = 1$, $u_1(x) u_2(x^{-1}) = 1$, d'où $u(x^{-1}) = u(x)^{-1}$ pour $u =
u_1 = u_2$, puis $u_3 = u$. Enfin $(*)$ avec $z = (xy)^{-1}$ donne $u(x)u(y) =
u((xy)^{-1})^{-1} = u(xy)$\footnote{La page fait le même calcul dans un ordre
légèrement différent (« $u_3(y) = u_2(y^{-1})^{-1}$, pareil $u_1(y) =
u_2(y^{-1})^{-1}$ ») ; il est correct. La phrase qui énonce les $u_i$ commence
au bas de la page 6 et finit en tête de la page 5 : les feuillets sont classés
dans le désordre.}.

C'est le théorème classique selon lequel les isotopies entre tables de groupes
normalisées en l'unité sont des isomorphismes ; le dossier 69 (page 116) en
fait le même calcul pour décrire les autotopies de la table d'un
groupe\footnote{Rattaché en général à Albert (1943) ; cité de mémoire, nom
d'« isotopie » de nous.}. La page 5 pose alors la question qui restera ouverte :
\emph{le foncteur $\varphi$ est-il une équivalence ?} --- c'est-à-dire : tout
carré latin vérifiant l'axiome I vient-il d'un système de bitorseurs ? La
réponse est non en général (il y a des quasigroupes qui ne sont isotopes à
aucun groupe) ; la condition qui caractérise l'image est aujourd'hui la
condition de Reidemeister, ou de clôture des quadrilatères\footnote{Noms et
réponse de nous ; les pages ne tranchent pas la question.}. Les pages 8 à 16
cherchent cette condition.

\pagerange{8}{12}
\section*{3. La loi de composition d'une relation pointée (pages 8 à 12)}

Partons d'un objet de $\mathcal{C}$ et d'un point $a = (a_1, a_2, a_3)$ de
$\Gamma$. L'axiome I fournit des applications $\alpha_1 : E_1 \to
\operatorname{Isom}(E_2, E_3)$, $\alpha_2 : E_2 \to \operatorname{Isom}(E_3,
E_1)$, $\alpha_3 : E_3 \to \operatorname{Isom}(E_1, E_2)$ ($\alpha_1(x)$
envoie $y$ sur l'unique $z$ avec $(x, y, z) \in \Gamma$, et ainsi de suite) ;
on a $\alpha_3(a_3) : a_1 \mapsto a_2$, etc. Posons
\[
\begin{aligned}
u_3 &= \alpha_2(a_2)\, \alpha_1(a_1) : E_2 \to E_1, \\
u_1 &= \alpha_3(a_3)\, \alpha_2(a_2) : E_3 \to E_2, \\
u_2 &= \alpha_1(a_1)\, \alpha_3(a_3) : E_1 \to E_3 ,
\end{aligned}
\]
et \emph{admettons} la condition
\[
u_1 u_2 u_3 = \mathrm{id} .
\]
Les $u_i$ forment alors un système transitif d'isomorphismes entre les $E_i$
(on a $u_3(a_2) = a_1$, etc.), qui permet de les identifier à un même
ensemble $G$ pointé par $e$. On a donc $\Gamma \subset G^3$ avec $(e, e, e)
\in \Gamma$, et, pour tous $x, y$, une unique solution en chaque place
(conditions a) et b) de la page 10). Soit $\sigma_1, \sigma_2, \sigma_3$ les
bijections de $G$ définies par
\[
(e, x, \sigma_1 x) \in \Gamma, \qquad (\sigma_2 x, e, x) \in \Gamma, \qquad
(x, \sigma_3 x, e) \in \Gamma .
\]
La page en tire $\sigma_2\sigma_1 = \sigma_3\sigma_2 = \sigma_1\sigma_3 =
\mathrm{id}$, donc $\sigma_1 = \sigma_2 = \sigma_3 = \sigma$ et $\sigma^2 =
\mathrm{id}$\footnote{Ces relations ne découlent pas des seuls axiomes a) et
b) : elles traduisent, après l'identification des $E_i$, l'hypothèse $u_1 u_2
u_3 = \mathrm{id}$ admise page 8, et c'est d'elle seule qu'elles viennent. Une
première relation, $(\sigma_3\sigma_2\sigma_1)^2 = \mathrm{id}$, est biffée ;
un premier énoncé b), barré de hachures, définissait déjà les $\sigma_i$. Le
début de la page 10 continue, selon toute apparence, la dernière ligne de la
page 8.}. On définit enfin une loi de composition par
\[
(x, y, z) \in \Gamma \iff z = \sigma(x \cdot y).
\]

\textbf{La structure obtenue} (page 12). Les conditions se traduisent ainsi :
les translations $z \mapsto zx$ et $z \mapsto xz$ sont bijectives ; $\sigma e =
e$ ; $e$ est unité à gauche et à droite ; $x \cdot \sigma x = e$, et comme
$\sigma^2 = \mathrm{id}$, $\sigma x \cdot x = e$. Autrement dit $(G, \cdot, e)$
est une \emph{boucle} (un quasigroupe à élément neutre) où tout élément a un
inverse bilatère $\sigma x$\footnote{Le mot « boucle » (en anglais
\emph{loop}) est de nous. La page a d'abord écrit « $\sigma x$ est l'inverse à
droite de $x$ », biffé puis rétabli sous une autre forme ; l'alinéa est marqué
d'un double trait en marge.}. La relation devient
\[
(x, y, z) \in \Gamma \iff (x \cdot y) \cdot z = e \iff z \cdot (x \cdot y) = e,
\]
et la page s'arrête sur la question qui reste : « $(xy)z = x(yz)$ ?? ».

\pagerange{13}{16}
\section*{4. L'associativité, et une identité de Moufang (pages 13 à 16)}

\textbf{Le cas d'un groupe} (page 13, « (A) »). Pour $E_1 = E_2 = E_3 = G$
un groupe, $\Gamma = \lbrace xyz = 1\rbrace$ et un point $(a, b, c)$ de $\Gamma$,
les applications $f_1$, $f_2$, $f_3$ définies par $(a, x, f_1(x)) \in \Gamma$,
$(f_2(x), b, x) \in \Gamma$, $(x, f_3(x), c) \in \Gamma$ valent
\[
f_1(x) = x^{-1}a^{-1}, \qquad f_2(x) = x^{-1}b^{-1}, \qquad f_3(x) =
x^{-1}c^{-1},
\]
leurs composés $g_1 = f_3 f_2 : x \mapsto b x c^{-1}$, $g_2 = f_1 f_3 : x \mapsto
c x a^{-1}$, $g_3 = f_2 f_1 : x \mapsto a x b^{-1}$, et
\[
g_1 g_2 g_3 (x) = (bca)\, x\, (cab)^{-1} = x
\]
puisque $abc = 1$ entraîne $bca = cab = 1$. C'est la condition admise page 8,
vérifiée pour un groupe ; on a aussi $g_1 g_2 g_3 = (f_3 f_2 f_1)^2$ et
$f_3f_2f_1(x) = b x^{-1} b$. L'involution $\sigma$ de la page 10 est alors $x
\mapsto x^{-1}$, et elle vérifie
\[
\sigma(xy) = \sigma y\, \sigma x,
\]
relation encadrée sur la page\footnote{En marge de la page 14, la même relation
est écrite $\sigma(xy) = \sigma x\, \sigma y$ ; dans un groupe non commutatif
c'est l'ordre de la page 13 qui est juste : l'inversion est un
anti-automorphisme.}.

\textbf{Sans associativité} (pages 14 et 15). La page 14 refait le calcul en
écrivant toutes les parenthèses, pour une boucle comme celle de la page 12. La
condition $g_1 g_2 g_3 = \mathrm{id}$ y devient une identité entre produits
parenthésés, que la page met sous la forme
\[
c\bigl((ax)b^{-1}\bigr) = \bigl(b^{-1}(xc)\bigr)a \qquad (abc = 1),
\]
puis, page 15, après plusieurs substitutions à deux encres, sous la forme
encadrée
\[
\bigl((bz)(xb)\bigr)c = b\bigl((zx)(bc)\bigr).
\]
Pour $z = 1$ elle se réduit à $\bigl(b(xb)\bigr)c = b\bigl(x(bc)\bigr)$, ce que
la page écrit aussi. C'est, à la flexibilité $b(xb) = (bx)b$ près, l'\emph{identité
de Moufang à gauche} $((bx)b)c = b(x(bc))$\footnote{Le rapprochement est de
nous. La page 15 est très surchargée ; la ligne « $z = 1$, $c = \ldots$ », dans
laquelle des lettres sont entourées, est d'une lecture d'ensemble douteuse, et
deux facteurs d'une formule intermédiaire sont illisibles. On ne reproduit pas
les calculs intermédiaires, dont la transcription donne l'état.}. Les boucles
de Moufang ne sont pas toutes des groupes (les octaves de Cayley de norme 1 en
donnent une) ; les pages ne disent pas si la condition de la page 8, appliquée
à tous les points de $\Gamma$, force l'associativité, et elles ne concluent
pas.

\textbf{Les petits ensembles} (page 16). La page, barrée de deux diagonales,
essaie les lois sur des ensembles à deux, trois, quatre éléments : sur $\lbrace
e, a\rbrace$ on trouve $aa = e$ ; sur $\lbrace e, a, b\rbrace$, $aa = b$, $ab =
ba = e$, $bb = a$ ; sur $\lbrace e, a, b, c\rbrace$ deux tables commencées, celle du
groupe de Klein ($aa = e$, $ab = c$, $ac = b$) et celle d'un groupe cyclique
($aa = b$, $ab = c$, $ac = e$). Ce sont bien les seules possibilités : toute
boucle d'ordre au plus 4 est un groupe, et la plus petite boucle non
associative a cinq éléments\footnote{Fait classique, de nous. Le bas de la page
revient aux relations $(ab)c = 1$, $(ax)y = 1$, \ldots, avec $f_1(x) =
\sigma(ax)$, $f_2(x) = (\sigma x)b^{-1}$ ; tête-bêche, à l'encre brune, quelques
formules reprennent l'associativité.}.

\pagerange{18}{22}
\section*{5. Un graphe sur un groupe, et un premier hexagone (pages 18 à 22)}

\textbf{Page 18.} Une construction nouvelle, sur un feuillet blanc, sans lien
écrit avec ce qui précède : un groupe $V$, une partie $I \subset V$, des
homomorphismes $\varphi_\lambda : V \to V_\lambda$ ($\lambda \in \Lambda$) avec
des parties $I_\lambda \subset V_\lambda$, un $V$-ensemble $Y$ et les
ensembles induits $X_\lambda = Y \wedge_V V_\lambda$, avec $t_\lambda : Y \to
X_\lambda$. Sur $S = \lbrace e\rbrace \amalg Y \amalg \coprod_\lambda
X_\lambda$ on met le graphe où $e$ est lié aux éléments des $X_\lambda$, $y \in
Y$ à $iy$ ($i \in I$), $x \in X_\lambda$ à $ix$ ($i \in I_\lambda$), et $y$ à
$t_\lambda(y)$\footnote{Au premier item le complément de « lié à tous les él.\
de » est surchargé et l'indice de $X$ illisible ; on lit les $X_\lambda$. Au
troisième la page porte « $i \in V_\lambda$ » ; on suit la transcription, qui
attend $I_\lambda$.}. « À quelles conditions est-ce là un graphe triadique ? »
Il annonce qu'il n'y a que les cas a) $V \simeq \mathbb{F}_2^{r+1}$, un
cardinal (illisible) valant $2^r$ avec $r \in \lbrace 0, 1, 2\rbrace$ ; b) $I$
est une base de $V$ ; et c), laissé vide. On reconnaît, sans que la page le
dise, la description d'un graphe de la famille autour d'un sommet $e$ : ses
voisins groupés par triangles, ses non-voisins $Y$ munis d'une action d'un
$\mathbb{F}_2$-espace vectoriel, et les trois valeurs $c = 2^r \in \lbrace 1,
2, 4\rbrace$ que les pages 120 et 127 retrouveront\footnote{Rapprochement de
nous ; la page n'en dit pas davantage et l'énoncé n'est pas démontré. La
description exacte autour d'un sommet est celle des pages 127 à 129 et 135
(section 24 ci-dessous) ; pour $c = 4$ les non-voisins forment un torseur sous un
espace de dimension 4 et l'ensemble $I$ des « translations voisines » n'y est
pas une base mais un système de cinq vecteurs de somme nulle.}.

\textbf{Pages 20 à 22.} Les vingt-sept apparaissent, dans la notation de
Schläfli : $\xi_1, \ldots, \xi_6$, les $\xi'_i$ et les quinze $\xi_{ij}$. La
page 20 dessine un hexagone dont les sommets sont des triplets ($\xi_1 \xi_2
\xi_3$, $\xi_4 \xi_5 \xi_6$, \ldots), « hexagone distingué ou triplets
typiques », et décompte, pour une telle configuration, les sommets liés à
aucun, un, deux des éléments fixés ($6$, $9$, $6$, \ldots) ; elle décrit aussi
l'ensemble comme fibré de degré 3 sur un ensemble $\varepsilon$ à deux
éléments, $E = \coprod_{\alpha \in \varepsilon} E_\alpha$, avec $E \times I$
fibré sur l'hexagone $I \times \varepsilon$\footnote{Les indices des triplets
aux sommets de l'hexagone ne se lisent pas tous ; les esquisses du bas de la
page (grilles $3 \times 3$, « $9\,6\,6\,6$ », « $9\,9\,9$ », « $40$ ») sont sans
légende.}. La page 21 note la correspondance « triades $\leftrightarrow$
hexagones », et les nombres $6^4 = 1296 = 2^4 3^4$ et $720 = 2^4 \cdot 3^2
\cdot 5$ ; la page 22 dessine les trois blocs de neuf
\[
\lbrace \xi_1, \xi_2, \xi_3, \xi'_1, \xi'_2, \xi'_3, \xi_{23}, \xi_{31},
\xi_{12}\rbrace, \quad \lbrace \xi_4, \ldots, \xi_{45}\rbrace, \quad \lbrace
\xi_{ij} \mid i \leq 3 < j\rbrace ,
\]
qui partagent les vingt-sept en trois pavés --- l'une des 40 décompositions
qu'on retrouvera page 25.

\pagerange{23}{30}
\section*{6. Dénombrements et racines (pages 23 à 30)}

La page 23 étudie un polyèdre (un octaèdre de sommets $p, p', q, q', r, r', s,
s'$ autour d'un point $a$) et la façon dont deux points sont liés ; la lecture
de la prose y est trop incertaine pour qu'on en tire un énoncé\footnote{On y
lit que $x, y$ « sont liés exactement par deux », l'une passant par $(r, a,
r')$, l'autre par $(s, a, s')$, et que deux points d'une partie $T'$ sont liés
si et seulement s'ils ont « des arêtes distinctes et convergentes à des sommets
distincts de $T$ » ; plusieurs mots illisibles.}. Les pages 25 à 29 sont des
dénombrements, en colonnes. Tous les nombres qui suivent sont justes :

\begin{itemize}
\item 27 sommets, 36 bibases, 40 « trihexagones », 45 triangles (tableaux des
  pages 28 et 29, où les bibases sont appelées « bitriades » puis
  « bibases »\footnote{À la page 28 le mot est « \emph{bitriades} », lu avec
  doute ; le tableau de la page 29 écrit « bibases ». Les « trois-neufs » de la
  page 29 sont les trihexagones.}) ;
\item 72 bases, $72 \cdot 6! / (72 \cdot 6) = 120$, $\binom{6}{3} = 20$ ;
\item 720 triades, « groupées en 120 hexagones qui se groupent par triples en
  40 trihexagones de triades » ; le stabilisateur d'une triade est
  $\mathfrak{S}_3 \times \mathfrak{S}_3 \times \mathfrak{S}_2$, d'ordre 72, et
  $51\,840/72 = 720$ ;
\item le stabilisateur d'un hexagone de triades est $(\mathfrak{S}_3 \times
  \mathfrak{S}_3 \times \mathfrak{S}_3) \cdot \mathfrak{S}_2$, d'ordre $6^3
  \cdot 2 = 432$, et il y en a $120 = 51\,840/432$ ; celui d'un trihexagone,
  décomposition des vingt-sept en trois blocs de neuf, est $\mathfrak{S}_3
  \cdot (\mathfrak{S}_3 \times \mathfrak{S}_3 \times \mathfrak{S}_3)$, d'ordre
  $1296$, et il y en a $40$ ;
\item $27 \cdot 16 \cdot 10 / 6 = 720$ (le dénominateur et le $16$ sont lus
  avec doute ; le calcul est juste : une droite, une droite disjointe, une
  troisième disjointe des deux).
\end{itemize}

Ces nombres sont les classiques : 72 « six », 36 doubles-six, 45 plans
tritangents, 120 paires de trièdres de Steiner, 40 triples de paires de
trièdres partageant les vingt-sept droites\footnote{Noms de nous. Les
vingt-sept dans la notation $\xi_{ij}$ apparaissent page 26 comme $E \times E'
\amalg E' \times E'' \amalg E'' \times E$ pour trois ensembles à trois éléments,
avec deux règles de liaison incomplètes ; c'est le modèle que les pages 51 et
52 mèneront à bien (section 11).}.

\textbf{Les racines} (pages 26 et 27). Avec $\eta$ la classe d'une droite du
plan, la page écrit $r_6 = \eta - \xi_1 - \xi_2 - \xi_3$ et $r_0 = 2\eta -
\xi$, et calcule $(2\eta - \xi) \cdot \xi_j$\footnote{La page écrit « $= 2$ ».
Avec la forme d'intersection ($\eta^2 = 1$, $\xi_i^2 = -1$, $\eta \cdot \xi_i =
0$) on trouve $(2\eta - \xi) \cdot \xi_j = 1$ ; la lecture de $\eta$ et $\xi$
est d'ailleurs incertaine à cet endroit.}. La page 27 dessine trois hexagones
dont les côtés portent des racines, et le diagramme de Dynkin de type $E_6$ :
$r_1 - r_2 - r_3 - r_4 - r_5$, avec $r_6$ attaché à $r_3$, où $r_i = \xi_i -
\xi_{i+1}$ ($i \leq 5$). Les racines des côtés sont $r_{i,j} = \xi_i - \xi_j$
(par exemple $r_{1,3} = r_1 + r_2$, $r_{4,6} = r_4 + r_5$), et
\[
r_{456} = \eta - \xi_4 - \xi_5 - \xi_6 = r_1 + 2r_2 + 3r_3 + 2r_4 + r_5 +
r_6 .
\]\footnote{La page écrit $r_{456} = 3(r_1 + r_2 + r_3) + 2r_4 + r_5 + r_6$, en
lisant avec doute le coefficient de $r_4$ et le signe devant $r_3$. Le calcul
direct donne $\xi_1 + \xi_2 + \xi_3 - \xi_4 - \xi_5 - \xi_6 = r_1 + 2r_2 + 3r_3
+ 2r_4 + r_5$ ; c'est la même correction que le dossier 69 fait à sa page 63.
On a alors $r_0 = r_6 + r_{456} = r_1 + 2r_2 + 3r_3 + 2r_4 + r_5 + 2r_6$, la
plus grande racine.}
Les six racines $r_1, r_2$ ; $r_4, r_5$ ; $r_0, r_6$ forment trois systèmes
$A_2$ deux à deux orthogonaux, et « engendrent un sous-groupe d'indice 3 » :
c'est vrai du \emph{réseau} qu'elles engendrent, d'indice 3 dans le réseau des
racines de $E_6$ (les discriminants valent $27$ et $3$) ; le sous-groupe de
$W(E_6)$ engendré par les six réflexions, $W(A_2)^3$ d'ordre 216, est
d'indice 240\footnote{La page écrit « dans $W$ », lu avec doute. Les trois
hexagones de la page 27 correspondent aux trois blocs de neuf de la page 22 :
la décomposition $A_2 \times A_2 \times A_2$ est un trihexagone, ce que le cadre
de la page 29 résume (« décomposition triple $A_2 \leftrightarrow$ triple bloc
de triades »).}.

\textbf{La proposition de la page 28.} Elle est trop fragmentaire pour être
énoncée en entier. Elle affirme que certaines parties ordonnées (« vrais
3-bisons », corrigé de « bitriades ») sont conjuguées sous le groupe, qu'une
telle partie est contenue dans exactement deux bases $I$, $I'$, et décrit
l'ensemble $K \cup K'$ de six sommets obtenu en retirant une partie $J$ ; dans
l'exemple, $K = \lbrace \xi_1, \xi_2, \xi_3\rbrace$, $K' = \lbrace \xi'_1,
\xi'_2, \xi'_3\rbrace$\footnote{Plusieurs mots en sont lus avec doute
(« bitriade », « respectivement », « intersections », « automorphisme ») et la
fin est illisible ; on n'en retient que ce que les exemples confirment.}. Le
bas de la page note, sans les développer, des pistes : $\mathfrak{S}_4$ et
$\mathbb{F}_2^4$, « quadriques affines de dim.~3 sur $\mathbb{F}_2$ », un
ensemble $A(i)$ à quatre éléments et un torseur $P(i)$ sous
$\mathbb{F}_2^{A(i)}$, une suite exacte $0 \to k \to V'(i) \to V(i) \to 0$, et
un « $2^3 \cdot 3 \cdot 5 = 120$ !!! ». La page 30 aligne les neuf paires de
triplets qui forment, trois par trois, les trois hexagones de la page 27.

\pagerange{32}{33}
\section*{7. Extensions centrales et la fonction $z_\rho$ (pages 32 et 33)}

Soit $V$ un groupe commutatif noté additivement, $I$ un ensemble à trois
éléments,
\[
\Gamma = \operatorname{Ker}(V^I \xrightarrow{\ \Sigma\ } V),
\]
$\omega$ l'ensemble des deux ordres circulaires sur $I$, torseur sous $\mathbb{Z}/2$,
et $V^\omega = V \wedge \omega$, $\mathbb{Z}/2$ opérant sur $V$ par le
signe\footnote{La page écrit d'abord $\mathbb{Z}$, corrigé en
$\mathbb{Z}/2$ (lecture incertaine des deux fois).}. Pour $i \in I$ on
définit $\varphi_i : V^\omega \to \Gamma$ par
\[
\varphi_i(x \wedge \rho)_i = 0, \qquad \varphi_i(x \wedge \rho)_{\rho i} = x,
\qquad \varphi_i(x \wedge \rho)_{\rho^2 i} = -x ,
\]
ce qui est bien défini puisque $(-x) \wedge \rho = x \wedge \bar\rho$ avec
$\bar\rho = \rho^{-1}$. On a $\sum_i \varphi_i = 0$ et la suite exacte
\[
0 \to V^\omega \xrightarrow{\ \mathrm{diag}\ } (V^\omega)^I \to \Gamma \to 0 .
\]
C'est, pour $V$ quelconque, le groupe $\Gamma = \lbrace (x, y, z) \mid x + y + z
= 0\rbrace$ des pages 1 à 16 dans le cas commutatif ; les $\varphi_i(V^\omega)$
sont les trois sous-groupes $\operatorname{Ker} p_i$.

On cherche les systèmes formés d'une extension centrale $0 \to Z \to
\widetilde\Gamma \xrightarrow{u} \Gamma \to 0$ et d'homomorphismes $\psi_i :
V^\omega \to \widetilde\Gamma$ relevant les $\varphi_i$. Fixons $\rho \in
\omega$ et $i_0 \in I$, $i_1 = \rho i_0$, $i_2 = \rho i_1$. Puisque $\sum
\varphi_i = 0$, le produit
\[
z(x, \rho, i_0) = \psi_{i_0}(x \wedge \rho)\, \psi_{i_1}(x \wedge \rho)\,
\psi_{i_2}(x \wedge \rho)
\]
est dans $Z$ ; comme $Z$ est central, il ne dépend pas de $i_0$ (une
permutation circulaire d'un produit égal à un élément central ne le change
pas). On note $z_\rho(x) = z(x, \rho)$, $x \in V$. En inversant le produit on
trouve
\[
z_{\bar\rho} = -z_\rho ,
\]
de sorte que la paire $(z_\rho)_{\rho \in \omega}$ est un élément de $Z^V
\wedge_{\pm 1} \omega$.

\textbf{Proposition} (page 33). Si $j = \rho i$, alors pour $x, y \in V$
\[
\bigl[\psi_i(x \wedge \rho), \psi_j(y \wedge \rho)\bigr] = z_\rho(x + y) -
z_\rho(x) - z_\rho(y),
\]
où $[g, h] = ghg^{-1}h^{-1}$\footnote{La page écrit le second membre avec le
signe opposé, $z_\rho(x) + z_\rho(y) - z_\rho(x + y)$, et le premier avec des
$\varphi$ que la transcription lit peut-être comme des $\psi$. On a refait le
calcul : en faisant passer $\psi_{i_0}(y)$ et $\psi_{i_1}(y)$ à travers
$\psi_{i_1}(x)\psi_{i_2}(x)$ dans $\psi_{i_0}(x+y)\psi_{i_1}(x+y)\psi_{i_2}(x+y)$,
les trois commutateurs qui apparaissent se réduisent à $[\psi_{i_0}(x),
\psi_{i_1}(y)]$, avec le signe indiqué ; on l'a contrôlé sur un modèle de type
Heisenberg. Le signe dépend de l'ordre des facteurs dans la définition de $z$,
pas de la convention de commutateur (les deux conventions usuelles coïncident
ici). Il est sans effet sur la suite, qui n'utilise que la bi-additivité.}.
En particulier le membre de droite, qui est le « défaut d'additivité » de
$z_\rho$, est bi-additif en $(x, y)$ (le commutateur dans une extension
centrale d'un groupe commutatif l'est) et symétrique (puisque $z_\rho(x+y) =
z_\rho(y+x)$).

La page 33 en tire, pour $\lambda = z_\rho$, les deux relations
\[
(*) \qquad \lambda(0) = 0, \qquad \lambda(x) + \lambda(y+u) + \lambda(y+v) +
\lambda(x+u+v) = \lambda(y) + \lambda(x+u) + \lambda(x+v) + \lambda(y+u+v)
\]
pour tous $x, y, u, v$\footnote{La seconde équation est très surchargée et
sa lecture d'ensemble incertaine ; c'est la forme que donne la page 35, où
elle est récrite, et c'est celle qui résulte de la proposition.}.

\pagerange{35}{39}
\section*{8. Applications quadratiques et le module universel (pages 35 à 39)}

La seconde relation $(*)$ dit que la différence seconde
\[
(\Delta_{u,v}\lambda)(x) = \lambda(x+u+v) - \lambda(x+u) - \lambda(x+v) +
\lambda(x)
\]
ne dépend pas de $x$ ; en la prenant en $x = 0$ elle vaut $\mu(u, v) =
\lambda(u+v) - \lambda(u) - \lambda(v)$, et la relation équivaut donc à :
\emph{$\mu$ est bi-additive}. Une application $\lambda : V \to Z$ avec
$\lambda(0) = 0$ et $\mu$ bi-additive est ce qu'on appelle une application
polynomiale de degré au plus 2 sans terme constant (ou application
« quadratique » au sens large)\footnote{Une ligne biffée la disait d'abord
« bilinéaire symétrique ».}.

Pour $\mu$ fixée, ces $\lambda$ forment un torseur sous le groupe des $f$
vérifiant $f(0) = 0$ et $\Delta_{u,v} f = 0$, c'est-à-dire sous les
homomorphismes $V \to Z$. La page ajoute que, pour $V$ libre de base $(e_i)$,
une forme bi-additive symétrique $\mu$ provient d'une $\lambda$ si et seulement
si les $\mu(e_i, e_i)$ sont dans $2Z$\footnote{La phrase, sur une page écrite au
crayon d'un trait pâle, est en partie illisible et s'arrête sur « \ldots ».
L'énoncé est juste pour un $\mathbb{Z}$-module libre (on prend $\lambda(e_i) =
\mu(e_i, e_i)/2$). Pour un module libre sur $\mathbb{Z}/n\mathbb{Z}$, comme la
page le pose, il faut en plus $\lambda(n e_i) = 0$, soit $n\lambda(e_i) +
\binom{n}{2}\mu(e_i, e_i) = 0$ ; pour $n = 2$, c'est $2\lambda(e_i) = -\mu(e_i,
e_i)$, la relation de la page 38.}.

\textbf{L'extension universelle} (page 37). Pour toute $z_\rho : V \to Z$
vérifiant les deux relations, il existe une extension $\widetilde\Gamma$ de
$\Gamma$ par $Z$ et des relèvements $\psi_i$ qui la redonnent, le tout
« déterminé à isomorphisme unique près »\footnote{Énoncé sans démonstration
sur la page. Une note marginale écrite verticalement, en partie illisible,
précise que l'extension universelle ne dépend pas, à isomorphisme canonique
près, du choix de $\rho$ : on y écrit $z_\rho(x) = \lambda(x \wedge \rho)$
avec $\lambda : V^\omega \to \mathfrak{Z}^\omega = \mathfrak{Z} \wedge_{\pm 1}
\omega$, et l'on vérifie $z_{\bar\rho} = -z_\rho$.}. Il y a donc un groupe
« universel » $\mathfrak{Z}_V$, quotient de $\mathbb{Z}^{(V)}$ par les
relations $(*)$, et une extension universelle de $\Gamma$ par
$\mathfrak{Z}_V$. La page 38 l'identifie au $\mathbb{Z}$-module universel qui
reçoit $V$ par une application polynomiale de degré au plus 2 sans terme
constant, qu'elle écrit $V + \Gamma^2_{\mathbb{Z}}(V)$\footnote{Ce module est
aujourd'hui noté $P_2(V)$ ; c'est le quotient $I_V / I_V^3$ de l'idéal
d'augmentation de l'anneau $\mathbb{Z}[V]$ (Passi). Le « $+$ » de la page doit
se lire comme une extension et non comme une somme directe : pour $V =
\mathbb{Z}/2$, $P_2(V) = \mathbb{Z}/4$ alors que $V \oplus \Gamma^2(V) =
\mathbb{Z}/2 \oplus \mathbb{Z}/4$. Noms de nous.}.

\textbf{Exemple} (page 38) : $V$ un espace vectoriel sur $\mathbb{F}_2$, de
base $I$\footnote{Ici $I$ désigne une base de $V$, et non plus l'ensemble à
trois éléments ; un premier essai, avec l'ensemble $D = V \smallsetminus
\lbrace 0\rbrace$ des droites, est barré et encadré.}. Toute $\lambda$ comme
plus haut est polynomiale : $\lambda(2x) = 2\lambda(x) + \mu(x, x)$ et $2x = 0$
donnent $\mu(x, x) = -2\lambda(x) \in 2Z$. D'où
\[
\mathfrak{Z}_V \simeq (\mathbb{Z}/4\mathbb{Z})^I \oplus
(\mathbb{Z}/2\mathbb{Z})^{\mathfrak{P}_2(I)} ,
\]
ce qui est juste : une telle $\lambda$ est déterminée par les $\lambda(e_i)$,
d'ordre divisant 4, et les $\mu(e_i, e_j)$, $i < j$, d'ordre divisant
2\footnote{Après $(\mathbb{Z}/4\mathbb{Z})^I$, un « $+\,\mathbb{Z}/4\mathbb{Z}$ »
est biffé.}.

\textbf{La forme canonique en dimension 2} (page 39). Si $V$ est de
dimension 2 sur $\mathbb{F}_2$, il y a sur $V$ une forme quadratique
canonique,
\[
q(x) = 0 \text{ si } x = 0, \qquad q(x) = 1 \text{ si } x \neq 0 ,
\]
c'est-à-dire la forme anisotrope $x^2 + xy + y^2$, dont la forme polaire est la
forme alternée non dégénérée de $V$. Elle donne une extension centrale
canonique de $\Gamma(V)$ par $\mathbb{Z}/2$ telle que
\[
\psi_{i_0}(x)\, \psi_{i_1}(x)\, \psi_{i_2}(x) = z \neq 0 \qquad (x \neq 0).
\]
C'est l'extension que le dossier 69 (pages 14 et 26) construit à partir du
groupe des quaternions $Q_8$ et dont il montre qu'elle n'existe, avec cette
propriété d'anisotropie, que pour $\dim V \leq 2$. Pour $\dim V = 2$, $\Gamma =
V^2$ a 16 éléments et $\widetilde\Gamma$ en a 32 : ce sont les cardinaux de
l'ensemble $\Gamma$ et du revêtement $\widetilde\Gamma$ des pages 101 à 104
pour le graphe cubique ($c = 4$)\footnote{Rapprochement de nous ; la page 39
s'arrête sur « (et un triplet $(i_0, i_1, i_2)$ sur $I$ \ldots) ».}.

\pagerange{41}{44}
\section*{9. Un brouillon sur les pentagones (pages 41 à 44)}

La moitié inférieure de la page 41, écrite tête-bêche, parle d'un
« dodécaèdre combinatoire » et de ses douze faces, vues comme les douze
structures de polygone d'un ensemble\footnote{Lectures incertaines
(« dodécaèdre », « icosaèdre », le symbole de l'ensemble).}. La page 43 est un
diagramme de flèches entre des ensembles construits sur un hexagone
(« Tritria », « Hex », « (tria$'$) », $S = \tau \times \varepsilon$,
$\widetilde S = \tau \times \varepsilon \times \omega(\tau)$\ldots), ébauche du
diagramme des revêtements de la page 72. La page 44 compte les pentagones sur
un ensemble à cinq éléments et les paires $\lbrace i, i+2\rbrace$ : $15 \cdot 6
\cdot 4 \cdot 2 \cdot 1 = 720 = 6!$, $10\,P = 720$, $P = 72$, $P_0 = 12$, ce
dernier étant le nombre des structures de pentagone sur un ensemble à cinq
éléments ($4!/2 = 12$).

\pagerange{46}{50}
\section*{10. Triades, hexagones et bitriangles dans le graphe cubique (pages 46 à 50)}

\textbf{Définitions} (page 46). Une partie d'un graphe est \emph{discrète} si
elle ne contient aucune arête ; les parties discrètes de cardinal 2, 3,
\ldots, 6 sont les doublets, triades, quadruplets, pentuplets, hexaplets. Un
graphe \emph{bi-discret} est un graphe $\Gamma = A \amalg A'$, réunion de deux
parties discrètes, tel que la relation « $x \in A$, $x' \in A'$ non liés » soit
une bijection de $A$ sur $A'$. La décomposition est unique, sauf si
$\operatorname{card} \Gamma = 4$ (le graphe est alors la somme de deux segments,
et une décomposition équivaut à une bijection entre ces segments) ; la
catégorie des graphes bi-discrets d'ordre $2n$ ($n \geq 3$) est équivalente à
$(\mathrm{Ens})_2 \times (\mathrm{Ens})_n$\footnote{La page écrit « d'ordre $n$
(card \ldots) », le cardinal étant illisible ; l'équivalence n'est juste que si
$n$ est le cardinal de chaque partie. On a ajouté l'hypothèse $n \geq 3$, qui
écarte le cas exceptionnel $n = 2$ que la page vient de signaler.}. Dans le
graphe cubique, les bases et bibases sont les hexaplets et les
bihexaplets ; une « bitriade » au sens de cette page est un graphe bi-discret à
six sommets, c'est-à-dire un hexagone (un 6-cycle), et la note marginale le dit
(« les bitriades sont exactement les hexagones »).

Deux triades $\tau$, $\tau'$ sont dites \emph{liées} si elles sont disjointes
et que leur réunion est discrète, c'est-à-dire une base. Deux parties
discrètes sont dites « W-liées » si elles sont disjointes et que leur réunion,
avec sa décomposition naturelle, est un graphe bi-discret\footnote{Le mot est
une capitale en forme de \emph{W} (ou \emph{V}) suivie de « -liées » ; « W-liées »
n'est qu'une approximation du tracé, gardée telle quelle par la transcription
et ici. La définition, telle qu'elle se lit, ne s'accorde pas avec l'usage qui
suit : deux triades dont la réunion est un hexagone sont, dans la proposition
ci-dessous, « biliées », et les triades « W-liées » sont les triades
\emph{opposées} (tout élément de l'une est lié à tout élément de l'autre),
comme le confirment les pages 49 et 53. On suit l'usage.}.

\textbf{Proposition} (pages 46 à 49). Soit $\Delta = (S, A)$ le graphe cubique
et $\mathrm{Tri}$ l'ensemble de ses 720 triades.

\begin{enumerate}
\item[(1)] Pour la relation « liées », $\mathrm{Tri}$ est une somme disjointe de
  120 hexagones.
\item[(2)--(3)] Deux triades distinctes d'un même hexagone sont liées
  (consécutives), W-liées (opposées) ou \emph{biliées} (à distance 2) ; deux
  triades sont biliées si et seulement si elles ne sont ni égales, ni liées,
  ni W-liées et qu'une troisième triade est liée aux deux\footnote{La page
  numérote deux points (3), et la phrase sur les triades « opposées » est en
  partie couverte de ratures ; on donne l'énoncé que les points (4) à (7)
  utilisent.}.
\item[(4)] Si $\tau'$, $\tau''$ sont les deux triades liées à $\tau$, alors
  $\tau' \amalg \tau'' = L(\tau)$, l'ensemble des sommets liés à aucun élément
  de $\tau$ ; c'est un hexagone (une bitriade), dont $\tau'$, $\tau''$ sont les
  deux parties.
\item[(5)] L'ensemble $\bar\tau = E(\tau)$ des sommets liés à tous les
  éléments de $\tau$ est une triade, l'\emph{opposée} de $\tau$.
\item[(6)] Les deux triades biliées à $\tau$ sont $\overline{\tau'} =
  L(\tau'') \smallsetminus \tau$ et $\overline{\tau''} = L(\tau')
  \smallsetminus \tau$, et l'hexagone des triades contenant $\tau$ est, dans
  l'ordre circulaire, $\tau, \tau'', \overline{\tau'}, \bar\tau,
  \overline{\tau''}, \tau'$.
\item[(7)] Pour la relation « biliées », $\mathrm{Tri}$ est une somme de 240
  triangles (« Tritria ») ; la relation « W-liés » entre triangles est une
  involution sans point fixe, dont les orbites correspondent aux hexagones.
\item[(8)] La réunion des triades d'un hexagone a 18 éléments, et son
  complémentaire (9 éléments) est un bitriangle ; l'application hexagone
  $\mapsto$ bitriangle ainsi obtenue est bijective.
\end{enumerate}

Dans le modèle de Schläfli, pour $\tau = \lbrace a_1, a_2, a_3\rbrace$ :
$L(\tau) = \lbrace a_4, a_5, a_6, c_{45}, c_{46}, c_{56}\rbrace$, hexagone
$a_4 - c_{45} - a_5 - c_{56} - a_6 - c_{46}$, de parties $\tau' = \lbrace a_4,
a_5, a_6\rbrace$ et $\tau'' = \lbrace c_{45}, c_{46}, c_{56}\rbrace$ ; $\bar\tau
= \lbrace b_4, b_5, b_6\rbrace$ ; et le bitriangle de (8) est $\lbrace c_{ij}
\mid i \leq 3 < j\rbrace$, à neuf droites, où $c_{ij}$ et $c_{kl}$ sont liés si
et seulement si $i \neq k$ et $j \neq l$\footnote{Exemple de nous, vérifié.
Tous les énoncés (1) à (8) le sont aussi sur le modèle.}.

\textbf{Le triangle de triades} (pages 48 à 50). (8) bis : si
$(\tau_i)_{i\in I}$ est un triangle de triades biliées et $x_i \in \tau_i$,
l'unique élément $x_j$ de $\tau_j$ non lié à $x_i$ et l'unique élément $x_k$ de
$\tau_k$ non lié à $x_i$ ne sont pas liés ; autrement dit la partie $\Gamma
\subset \prod \tau_i$ des triplets deux à deux non liés a ses trois projections
bijectives --- c'est l'axiome I de la page 1, et un système transitif
d'isomorphismes entre les $\tau_i$. (9) La réunion $X$ des trois triades est un
bitriangle, les $\tau_i$ sont l'une de ses deux familles de triades disjointes,
et $X \to I \times J$ est bijectif, $J$ indexant l'autre famille. (10) Se
donner un triangle de triades biliées équivaut à se donner un bitriangle et
le choix de l'une de ses deux familles de triades (une « coorientation »).
(11) Deux triades sont biliées si et seulement si elles sont disjointes et
contenues dans un même bitriangle, et celui-ci est unique. (12) Le
complémentaire $\widetilde X = S \smallsetminus X \smallsetminus \overline X$
du bitriangle et de son « W-lié » est encore un bitriangle ; le groupe
$\widetilde G$ engendré par les symétries $\sigma_t$ attachées aux triangles $t$
de $\widetilde X$ (voir la section 21) s'envoie isomorphiquement sur le groupe
des automorphismes « intérieurs » de $\widetilde X$, et ses orbites sur $S
\smallsetminus \widetilde X$ sont trois ensembles à six éléments de la forme
$\tau \cup \bar\tau$ ; le sous-groupe des produits pairs a pour orbites les six
triades de l'hexagone\footnote{Les points (12) b) et c) sont écrits dans une
prose surchargée (« bihexaH », « $\mathfrak{S}_{\mathbb{F}} \times
\mathfrak{S}_{\mathbb{Q}}$ » pour les deux familles de triangles de $\widetilde
X$) ; on en donne ce qui se lit sûrement. La numérotation (11) revient deux
fois, pages 49 et 50 ; c'est celle du feuillet.}.

\pagerange{50}{53}
\section*{11. Trois bitriangles disjoints, et deux modèles (pages 50 à 53)}

\textbf{Proposition} (pages 50 et 51). Soit $(X_\alpha)_{\alpha \in \Lambda}$
une famille de trois bitriangles deux à deux disjoints (ils partagent donc les
vingt-sept sommets). Pour chaque $\alpha$, le groupe des automorphismes
intérieurs de $X_\alpha$ a trois orbites sur $S \smallsetminus X_\alpha$ ;
leur ensemble $I_\alpha$ a trois éléments, et l'on a des applications $p_{\alpha
\beta} : X_\beta \to I_\alpha$ ($\beta \neq \alpha$) dont les fibres sont l'une
des deux familles de triades de $X_\beta$. Alors :

\begin{enumerate}
\item[a)] pour $\Lambda = \lbrace \alpha, \beta, \gamma\rbrace$, les deux familles
  de fibres de $X_\alpha \to I_\beta$ et $X_\alpha \to I_\gamma$ sont distinctes,
  et $X_\alpha \xrightarrow{\sim} I_\beta \times I_\gamma$ ; d'où une bijection
  naturelle
  \[
  S \simeq \coprod_{\alpha \in \Lambda} \prod_{\beta \neq \alpha} I_\beta =
  (I_2 \times I_3) \amalg (I_3 \times I_1) \amalg (I_1 \times I_2) ;
  \]
\item[c)] deux sommets distincts de $X_\alpha$ sont non liés si et seulement
  s'ils ont une coordonnée commune ;
\item[d)] $x_\alpha \in X_\alpha$ et $y_\beta \in X_\beta$ ($\alpha \neq \beta$)
  sont liés si et seulement s'ils ont même coordonnée dans $I_\gamma$, $\gamma$
  le troisième indice.
\end{enumerate}

\textbf{Un modèle} (page 52). De façon équivalente, soit $p : I \to \Lambda$
une application d'un ensemble à neuf éléments sur un ensemble à trois, à fibres
$I_\alpha$ de cardinal 3. Les sommets sont les parties $A \subset I$ de
cardinal 2 sur lesquelles $p$ est injective (27 parties), $X_\alpha$ étant
formé de celles pour lesquelles $p(A) = \Lambda \smallsetminus \lbrace
\alpha\rbrace$ ; deux sommets $A \neq B$ sont liés si et seulement si
\[
\text{ou bien } p(A) = p(B) \text{ et } A \cap B = \emptyset, \qquad \text{ou
bien } p(A) \neq p(B) \text{ et } A \cap B \neq \emptyset .
\]
Les deux descriptions donnent bien le graphe cubique\footnote{Vérifié : dans
les deux cas le graphe est fortement régulier de paramètres $(27, 10, 1, 5)$, ce
qui le caractérise (dossier 69, section 1). C'est un modèle voisin de celui de
la page 13 du dossier 69, construit sur deux ensembles à trois éléments.}.

\textbf{(12)} La catégorie des graphes cubiques munis d'une décomposition en
trois bitriangles disjoints est équivalente à celle des ensembles à neuf
éléments munis d'une partition de type $(3, 3, 3)$, et le groupe
d'automorphismes d'une telle structure est
\[
\mathfrak{S}_3 \cdot (\mathfrak{S}_3 \times \mathfrak{S}_3 \times
\mathfrak{S}_3), \qquad \text{d'ordre } 6^4 = 1296 = 2^4 \cdot 3^4 .
\]\footnote{La page
écrit « d'ordre $3 \cdot 6^4$ », lu avec doute ; c'est $6 \cdot 6^3 = 6^4$. La
note marginale porte « $3^4 2^4$ ».}
Comme $|W(E_6)| = 2^7 \cdot 3^4 \cdot 5$, ce groupe contient un 3-sous-groupe
de Sylow de $\operatorname{Aut}(\Delta)$, ce que la page exprime par « c'est
donc presque un sous-groupe de Sylow »\footnote{La note marginale qui
développe ce point, sur le normalisateur, n'est lue que par fragments.}. La
catégorie des couples $(\Delta, X)$, $X$ un bitriangle, est équivalente à
$(\mathrm{Ens})_3 \times (\mathrm{Ens})_{3,3}$ : l'ensemble des trois orbites
associé à $X$, et l'ensemble des six sommets d'une famille de triades de $X$
avec sa partition naturelle (ordre $6 \cdot 72 = 432$, et $51\,840 / 432 =
120$).

\textbf{Triades et hexagones dans le modèle} (page 53). Les 18 triades
contenues dans l'un des $X_\alpha$ correspondent aux couples $(i, \alpha) \in I
\times \Lambda$ avec $p(i) \neq \alpha$ : à $(i, \alpha)$ la triade
$\tau_{i,\alpha}$ des $A$ avec $p(A) = \Lambda \smallsetminus \lbrace
\alpha\rbrace$ et $A \ni i$. La triade W-liée à $\tau_{i\alpha}$ est
$\tau_{i\beta}$, $\beta$ l'autre élément de $\Lambda \smallsetminus \lbrace
p(i)\rbrace$ ; les deux triades liées sont les $\tau_{j\beta}$, $\tau_{k\beta}$,
$j, k$ les deux autres éléments de la fibre de $i$. Ces 18 triades forment
trois hexagones, en bijection avec $\Lambda$ (à l'hexagone correspond le
complémentaire de sa réunion, l'un des $X_\alpha$).

\pagerange{53}{60}
\section*{12. Triangles de triades, la tour, et les types de triades (pages 53 à 60)}

(14) Les triangles de triades biliées correspondent aux repères $(\alpha,
\beta, \gamma)$ de $\Lambda$, l'hexagone $H \simeq I_\gamma \times \Lambda_\gamma$
($\Lambda_\gamma = \lbrace \alpha, \beta\rbrace$), $\tau_{i\alpha}$ et
$\tau_{j\beta}$ étant liées si et seulement si $i \neq j$ et $\alpha \neq
\beta$. Les deux opérations « passer au triangle W-lié » ($-$) et « passer à
l'autre triangle du même bitriangle » ($\sim$) vérifient $\overline{T_{\alpha
\beta\gamma}} = T_{\beta\alpha\gamma}$ et $\widetilde{T_{\alpha\beta\gamma}} =
T_{\alpha\gamma\beta}$, et engendrent sur les triangles de triades un groupe
isomorphe à $\mathfrak{S}_3$ (page 56). (15) Les triangles de $X_\alpha$
correspondent aux bijections $I_\beta \to I_\gamma$, et la symétrie associée
agit sur $\Lambda$ en fixant $\alpha$ et en échangeant $\beta$, $\gamma$
(page 57)\footnote{Le milieu de la page 57 n'est pas lu de près par la
transcription ; on en donne le début et la fin.}.

\textbf{La tour} (pages 57 à 59). (16)--(17) Pour un bitriangle $X$, les
sous-graphes « de type $\mathcal{T}_1$ » contenant $X$ sont les $X \cup
\mathcal{X}_i$, $\mathcal{X}_i$ l'une des trois orbites de $S \smallsetminus X$ ;
pour un $\mathcal{T}_1$ $\mathcal{X}$, le complémentaire $S \smallsetminus
\mathcal{X}$ est une bibase, d'où une bijection entre les $\mathcal{T}_1$ du
graphe cubique et ses 36 bibases. Les systèmes $S \supset \mathcal{X}
\supset X$ forment une catégorie équivalente à $(\mathrm{Ens})_{3,3} \times
(\mathrm{Ens})_2$, et les systèmes $S \supset \mathcal{X} \supset X \supset t$
($t$ une triade de $X$) à $(\mathrm{Ens})_3 \times (\mathrm{Ens})_2 \times
(\mathrm{Ens})_2$\footnote{La page 58 décrit les foncteurs (« l'ensemble des
points de l'une des triades de $X$ par la relation d'équivalence définie par
$t$, l'ensemble des deux coorientations, l'ensemble des deux orbites ») ; on ne
les reprend pas en détail. Le dernier membre de la phrase de (17) est
illisible.}.

\textbf{Le graphe à quinze sommets} (page 59). Pour un ensemble $B$ à six
éléments, soit $\mathcal{T}_1(B)$ le graphe de sommets $\mathfrak{P}_2(B)$, deux
paires étant liées si et seulement si elles sont disjointes. C'est un
$\mathcal{T}_1$, et $B \mapsto \mathcal{T}_1(B)$ est une équivalence de
$(\mathrm{Ens})_6$ sur la catégorie des $\mathcal{T}_1$ ; un quasi-inverse
associe à $\mathcal{X}$ l'ensemble de ses six parties discrètes à cinq
éléments, chaque sommet étant dans exactement deux d'entre elles\footnote{On l'a
vérifié : les parties discrètes à cinq éléments de $\mathcal{T}_1(B)$ sont les
six « étoiles » $\lbrace \lbrace i, j\rbrace \mid j \neq i\rbrace$. Les triangles de
$\mathcal{T}_1(B)$ sont les quinze partitions de $B$ en trois paires (les
« synthèmes » de Sylvester) ; le graphe est le quadrangle généralisé $W(2)$,
c'est-à-dire la configuration de Cremona--Richmond. Noms de nous.}. Un
bitriangle dans $\mathcal{X}$ équivaut à une partition de type $(3, 3)$ de $B$ :
ses neuf sommets sont les paires qui rencontrent les deux parties.

\textbf{Les types de triades} (page 60). Soit $\Delta$ muni d'une partition
$S = \coprod_\alpha X_\alpha$ en trois bitriangles, et $G$ son groupe
d'automorphismes (d'ordre 1296). $G$ a quatre orbites sur les triades : celle
des 18 triades contenues dans un $X_\alpha$, et trois autres, de types I, II,
III, que la page décrit par des dessins (trois colonnes de trois points)\footnote{La
page avait écrit « trois orbites », corrigé en « quatre » ; la correction est
juste. On a calculé les cardinaux : $18$, $162$, $216$, $324$ (total 720) ; on ne
peut attribuer sûrement les trois derniers aux types I, II, III, les dessins
n'étant pas reproduits.}. Les triades contiguës à une triade de type I sont l'une
de type I, l'autre de type II, de sorte qu'un hexagone de ce genre est de type
(I, II, II, I, II, II) en partant d'un sommet de type I ; la phrase sur le type
III se poursuit au-delà de la page 60, sur des feuillets qui ne la reprennent
pas.

\pagerange{61}{66}
\section*{13. Les 120 pavés vus de l'un d'eux (pages 61 à 66)}

La page 61 dessine un hexagone de triades qui n'est contenu dans aucun des
$X_\alpha$ (de deux types I et quatre types II, ou l'inverse)\footnote{Les
étiquettes des sommets de l'hexagone dessiné sont en grande partie douteuses.}
et note que tout triangle de triades biliées est formé, soit de triades d'un
même $X_\alpha$, soit de deux triades de type I et une de type II, soit de
triades de type III. Page 62 : les bitriangles correspondent aux couples formés
d'un sommet $a$ et d'une bijection convenable, et, pour un bitriangle $X$
donné, les 120 bitriangles se répartissent ainsi :

\begin{itemize}
\item $X$ lui-même : 1 ;
\item les deux bitriangles disjoints de $X$ : 2 ;
\item ceux qui coupent $X$ suivant deux triangles de même sommet (5 sommets) :
  27 ;
\item ceux qui coupent $X$ suivant un doublet : 54 ;
\item ceux qui coupent $X$ suivant un triangle : 36 ;
\end{itemize}
soit 120 au total.

C'est exact\footnote{Vérifié. La page a d'abord écrit « $3 \times 10$ », corrigé en
« $3 \times 18$ », pour les 54, et « $6 \times 6$ » pour les 36 ; le premier essai,
entre crochets, est barré.}. Page 63 : ceux qui coupent $X$ suivant un
triangle dans chaque $X_\alpha$ correspondent à un système transitif
d'isomorphismes entre les $I_\alpha$, six pour un triangle donné. Il note un
programme : étudier la position relative de deux $\mathcal{T}_1$, puis d'une
base et d'un $\mathcal{T}_1$ ; il n'est pas suivi.

\textbf{Triades et hexagones} (pages 63 à 65). (20) Une triade $\tau$ est
contenue dans un unique bitriangle, formé de $\tau$ et des six sommets liés à
exactement deux éléments de $\tau$ ; son complémentaire dans ce bitriangle est
un hexagone, d'où une bijection triades $\to$ hexagones\footnote{« Exactement » est
ajouté en interligne ; sans lui l'énoncé serait faux (les trois sommets liés
aux trois éléments de $\tau$ n'y sont pas). Dans le modèle, pour $\tau =
\lbrace a_1, a_2, a_3\rbrace$, ce bitriangle est $\tau \cup \lbrace b_1, b_2,
b_3\rbrace \cup \lbrace c_{12}, c_{13}, c_{23}\rbrace$.}. La bijection inverse associe à un hexagone $Y$ les sommets hors
de $Y$ liés à exactement quatre de ses sommets, deux à deux opposés, et qui
correspondent ainsi aux trois diagonales. Plus finement, pour une triade $\tau$
et $x \notin \tau$, soit $\nu_\tau(x)$ le nombre d'éléments de $\tau$ liés à $x$ ;
on a
\[
\begin{array}{lll}
\nu_\tau(x) = 0 : & 6 \text{ sommets,} & \text{un hexagone } h(\tau) ; \\
\nu_\tau(x) = 1 : & 9 \text{ sommets,} & \text{un bitriangle } \operatorname{bit}(\tau) ; \\
\nu_\tau(x) = 2 : & 6 \text{ sommets,} & \text{un hexagone } \bar h(\tau),
\text{ « opposé » à } h(\tau) ; \\
\nu_\tau(x) = 3 : & 3 \text{ sommets,} & \text{la triade opposée } \bar\tau .
\end{array}
\]
(page 64), et $\operatorname{bit}(\tau) = \operatorname{bit}(\bar\tau) =
\operatorname{bit}(\tau') = \operatorname{bit}(\tau'') = \operatorname{bit}(\bar\tau')
= \operatorname{bit}(\bar\tau'')$ pour les six triades de l'hexagone de triades
de $\tau$ (page 65)\footnote{Vérifié sur le modèle. Le symbole lu
$\mathrm{\acute Et}(x)$ (les éléments de $\tau$ liés à $x$) est douteux ; la marge de
la page 65 porte « détailler tout ça dans énoncés en forme ».}.

\textbf{Comptes} (page 66). La page vérifie $27 \cdot 16 \cdot 10 = 6 \cdot 720$
(les triades ordonnées) et $45 \times 16 = 720$, nombre des bitriangles munis
d'un triangle (chaque triangle est dans 16 bitriangles) et nombre des triades.
Elle note les équivalences « $A_4 \Leftrightarrow A_5 \Leftrightarrow P_6
\Leftrightarrow$ bitriangle épinglé », reprises page 133\footnote{Deux autres
nombres de la page, « $2^7 \cdot 3^3 \cdot 5 =$ Nb de hexagones épinglés » et
« $2^6 \cdot 3 \cdot 5 =$ Nb de hexagones $=$ Nb de bitriades », ne
s'accordent pas avec les comptes : il y a 720 hexagones et $720 \cdot 12 =
2^6 \cdot 3^3 \cdot 5$ hexagones épinglés (munis d'un sommet et d'un sens). On
ne les retient pas.}.

\pagerange{68}{70}
\section*{14. Le pavé et son complémentaire (pages 68 à 70)}

À partir de la page 68 le bitriangle s'appelle \emph{pavé} (« anciennement :
bitriangle »), et le graphe est noté $\Delta = (S, A, T)$. Soit $X \subset S$
un pavé.

\begin{enumerate}
\item[1)] Pour $x \in S \smallsetminus X$, l'« ombre » $\operatorname{Omb}_X(x)$
  des sommets de $X$ liés à $x$ est une triade de $X$ ; $\varphi_X : S
  \smallsetminus X \to \operatorname{Tria}(X)$ est surjective.
\item[2)] La fibre de $\varphi_X$ en une triade $t$ est une triade $\bar{\bar
  t}$, l'unique triade disjointe de $t$ dont tous les éléments sont liés à tous
  ceux de $t$ (l'opposée de $t$).
\item[3)] Les deux triades $t'$, $t''$ de $X$ parallèles à $t$ (de la même
  famille) ont pour réunion l'ensemble des sommets hors de $t$ liés à exactement
  deux éléments de $t$\footnote{La page parle ici de triades « liées » à $t$, au
  sens nouveau ; le texte entre parenthèses est en partie biffé et l'ancien sens
  du mot (page 46) s'y mêle. On donne l'énoncé que vérifie le modèle.}.
\item[4)] Les six triades $\tau, \tau', \tau''$ et $\bar\tau, \bar\tau',
  \bar\tau''$ forment un hexagone de triades.
\item[5)] Si $\rho, \rho', \rho''$ sont les trois triades de l'autre famille de
  $X$, les opposées $\bar\rho$, \ldots\ forment un second pavé $X_2$ ; $X_1$
  (formé des $\bar\tau$) et $X_2$ sont les seuls pavés contenus dans $S
  \smallsetminus X$, ils le partagent, et ils sont en bijection avec l'ensemble
  $\varepsilon$ des deux familles de triades de $X$.
\item[6)--7)] Les triades restantes de $X_1$ et $X_2$ se correspondent par
  opposition et forment encore un hexagone, dont l'ensemble $I_0$ des trois
  diagonales donne des bijections $X_1 \simeq I_1 \times I_0$, $X_2 \simeq I_2
  \times I_0$ ; pour $i \in I_0$, $X \cup \pi^{-1}(i)$ est un $\mathcal{T}_1$,
  et l'on obtient ainsi les trois $\mathcal{T}_1$ contenant $X$.
\item[8)] Ces fibres sont aussi les orbites, sur $S \smallsetminus X$, du
  groupe engendré par les symétries des triangles de $X$.
\end{enumerate}

C'est la décomposition de la page 51 vue depuis l'un des trois pavés\footnote{La
fin de 6), qui se poursuit en tête de la page 70, est en partie illisible.}.

\pagerange{72}{77}
\section*{15. Le diagramme des revêtements (pages 72 à 77)}

La page 72 dresse la liste des ensembles en jeu, avec leurs cardinaux, et un
diagramme d'applications dont les fibres ont toutes deux ou trois éléments :
\[
\begin{array}{ll}
\mathrm{Hex} & \text{hexagones de triades (« grands hexagones ») } \simeq
\text{ pavés : } 120 ; \\
\mathrm{hex} & \text{hexagones de sommets (« petits ») } \simeq \text{ triades
} (\mathrm{tria}) : 720 ; \\
\mathrm{triHex} & \text{triples distingués de grands hexagones : } 40 ; \\
\mathrm{Tria}' & \text{paires de triades opposées (diagonales des grands
hexagones) : } 360 ; \\
\mathrm{Trtria} & \text{triangles de triades biliées } \simeq \text{ pavés
orientés : } 240 .
\end{array}
\]
\begin{tikzcd}
\mathrm{Tria}' \arrow[d, "3"'] & \mathrm{tria} \simeq \mathrm{hex} \arrow[l, "2"'] \arrow[d, "3"] \\
\mathrm{Hex} \arrow[d, "3"'] & \mathrm{Trtria} \arrow[l, "2"] \\
\mathrm{triHex} &
\end{tikzcd}

Le carré est cartésien\footnote{La page porte « 720 » en regard de
$\mathrm{Hex}$ ; c'est 120, comme le donnent les degrés du diagramme ($360/3$)
et la page 46. La mention « (cart) » est au milieu du carré ; au-dessus de la
flèche supérieure, un petit mot « bib » lu avec doute.}. Les revêtements de
degré 2 sont les passages à la triade opposée et au triangle opposé ; et tout
le diagramme se déduit de $\mathrm{Tria}' \xrightarrow{3} \mathrm{Hex}
\xrightarrow{3} \mathrm{triHex}$, $\mathrm{Trtria}$ étant le revêtement double
des orientations associé au revêtement de degré 3 $\mathrm{Hex} \to
\mathrm{triHex}$.

\textbf{L'hexagone combinatoire} (pages 73 à 76). Un petit hexagone $h$ est
formé de deux triades parallèles $t = \lbrace a, b, c\rbrace$, $t' = \lbrace a',
b', c'\rbrace$ (sommets $a, b', c, a', b, c'$ dans l'ordre) ; il y a un unique
pavé $P \supset h$, et $t'' = P \smallsetminus h$ est la troisième triade
parallèle, son élément $a''$ correspondant à la diagonale $\lbrace a,
a'\rbrace$. Un hexagone abstrait équivaut à un couple $(\tau, \varepsilon)$
d'un ensemble à trois éléments (ses diagonales) et d'un ensemble à deux
éléments (ses deux triangles inscrits), de sommets $\tau \times \varepsilon$,
$(a, \alpha)$ et $(b, \beta)$ étant voisins si et seulement si $a \neq b$ et
$\alpha \neq \beta$ ; ses orientations sont celles de $\tau$. Pour une triade
$t$ décrite dans le modèle de la page 52 par $(i, \beta)$, $p(i) = \alpha$, le
petit hexagone correspond à $(I_\gamma, I_\alpha \smallsetminus \lbrace
i\rbrace)$ et le grand à $(I_\alpha, \lbrace 0, 1\rbrace)$, et la page note qu'il
n'y a pas de relation universelle entre leurs deux ensembles d'orientations.

La « présentation symétrique » de la page 75 : un hexagone détermine
l'ensemble $\tau$ de ses diagonales (canoniquement isomorphe à celui de ses
codiagonales, paires de côtés opposés), l'ensemble $\varepsilon$ de ses triangles
inscrits, l'ensemble $\varepsilon'$ de ses triangles « coinscrits » (triples de
côtés deux à deux non adjacents), et un isomorphisme $\varphi : \varepsilon
\wedge \varepsilon' \simeq \omega(\tau)$ ; sommets $\tau \times \varepsilon$,
côtés $\tau \times \varepsilon'$, et $(s, \alpha)$ incident à $(t, \beta)$ si et
seulement si $t = \varphi(\alpha \wedge \beta)$\footnote{Formulation de la page,
le premier terme de la règle d'incidence étant surchargé. Le « $1_{\mathcal{T}_2}$ »
qu'elle écrit pour l'élément trivial est peut-être un $1_{\mathbb{Z}/2}$.}. Le
passage à l'hexagone dual échange $\varepsilon$ et $\varepsilon'$ et change
$\varphi$ en $-\varphi \circ \sigma$. Le groupe des automorphismes est le
sous-groupe de $\operatorname{Aut}(\tau) \times \operatorname{Aut}(\varepsilon)
\times \operatorname{Aut}(\varepsilon')$ formé des $(u, \alpha, \beta)$ avec
$\alpha\beta\,\mathrm{sg}(u) = 1$, d'ordre 12 : c'est $D_6$.

Les pages 76 et 77 dessinent le « cube » des revêtements entre ensembles
attachés à un hexagone (sommets, côtés, diagonales, triangles inscrits et
coinscrits, orientations, repères), puis le grand diagramme reliant les
ensembles de configurations du graphe cubique, avec leurs groupes
($\mathfrak{S}_3 \times \mathfrak{S}_2$ pour les repères d'un grand hexagone,
etc.) ; la transcription les décrit et ne les restitue pas, et l'on fait de
même. La page 77 note enfin les sous-groupes de $\mathfrak{S}_3 \times
\mathfrak{S}_2$ : trois sous-groupes d'indice 2, trois d'indice 3
(conjugués), et des sous-groupes d'indice 6 répartis en trois classes,
dénombrement que la page 80 précise\footnote{Les sous-groupes d'ordre 2 sont
sept, en trois classes de conjugaison (le centre, et deux classes de trois) ;
c'est ce que dit la page 80.}.

\pagerange{78}{80}
\section*{16. Les polygones de trois à six côtés (pages 78 à 80)}

\begin{enumerate}
\item[1)] $n = 3$ : $\mathrm{Pol}_3 \simeq (\mathrm{Ens})_3$, et la dualité
  (sommets $\leftrightarrow$ côtés) est canoniquement isomorphe à l'identité,
  un côté correspondant au sommet opposé ; de même pour tout $n$ impair. Groupe
  $\mathfrak{S}_3$.
\item[2)] $n = 4$ : $\mathrm{Pol}_4 \simeq (\mathrm{Ens})_{2,2}$ (un carré est
  l'ensemble de ses sommets avec sa partition en diagonales), et plus généralement
  $\mathrm{Cube}_d \simeq (\mathrm{Ens})_{2, \ldots, 2}$ ($d$ fois). L'ensemble
  $C_r$ des repères d'un carré est un torseur sous $D_4$ (d'ordre 8) ; les
  ensembles du diagramme (sommets, côtés, diagonales, codiagonales, repères
  modulo la symétrie centrale, orientations) en sont des quotients, et la page
  constate qu'ils représentent chacune des huit classes de conjugaison de
  sous-groupes de $D_4$ exactement une fois\footnote{Huit classes : le groupe
  trivial, le centre, deux classes de symétries non centrales, le sous-groupe
  cyclique d'ordre 4, deux sous-groupes de Klein, $D_4$. La phrase qui le dit
  est en partie illisible, et l'attribution de quelques flèches du diagramme
  autour de $C_\omega$ et $C_{d^*}$ est douteuse.}.
\item[3)] $n = 6$ : $P_6 \simeq P_3 \times (\mathrm{Ens})_2 \simeq
  (\mathrm{Ens})_3 \times (\mathrm{Ens})_2$ (et de même pour $n = 2m$, $m$
  impair) : c'est l'hexagone $t \times \omega$ de la page 74. Les diagonales et
  les codiagonales sont canoniquement en bijection ; on a $H_\omega \simeq
  \omega(H_d) \simeq H_{tr} \wedge H_{tr^*}$, et les ensembles attachés à
  l'hexagone sont les espaces homogènes de
  \[
  D_6 \simeq \mathfrak{S}_2 \cdot \mathbb{Z}/6\mathbb{Z} \simeq \mathfrak{S}_2
  \times \mathfrak{S}_3
  \]
  (d'ordre 12) correspondant à ses classes de sous-groupes : trois sous-groupes
  conjugués d'indice 3 (stabilisateurs des diagonales), trois sous-groupes
  d'indice 2 non conjugués (stabilisateurs des points de $H_{tr}$,
  $H_{tr^*}$, $H_\omega$), un sous-groupe d'indice 4, et trois classes de
  sous-groupes d'ordre 2 (le centre, stabilisateur des points de $H_{p'}$, et
  deux classes de trois, stabilisateurs des sommets et des côtés).
\end{enumerate}

La page 80 s'arrête sur l'observation que les automorphismes de ces espaces
homogènes sont, eux aussi, lisibles sur le diagramme. C'est la
même manière de faire que pour les polygones du dossier 75 : un objet
combinatoire est remplacé par la liste des ensembles qu'il détermine, et ces
ensembles par les sous-groupes de son groupe d'automorphismes.

\pagerange{81}{85}
\section*{17. Carrés et hexagones dans un pavé (pages 81 à 85)}

Soit $X$ un pavé ($3 \times 3$). Un \emph{carré} de $X$ est un 4-cycle induit :
deux lignes et deux colonnes. (I) Les carrés de $X$ sont en bijection avec les
sommets de $X$, à un carré $C$ correspondant son \emph{centre} $a$, l'unique
sommet de $X$ lié à aucun sommet de $C$ (à l'intersection de la ligne et de la
colonne restantes). Deux sommets non liés sont dans une unique triade ; deux
paires de sommets non liés forment un carré si et seulement si elles ont même
« centre », et l'une détermine l'autre : la donnée d'un carré à diagonale
marquée équivaut à celle d'une paire non liée, ou d'une triade pointée. D'où
les isomorphismes
\[
\mathrm{Car}_d(X) \simeq \mathrm{tria\,pt}(X), \qquad \mathrm{Car}_s(X) \simeq
\mathrm{tria\,ép}(X),
\]
entre carrés à diagonale (resp.\ à sommet) marqué(e) et triades pointées (resp.\
épinglées)\footnote{Les indices $d$, $s$ sont lus d'après la page suivante.}.
Les quatre sommets de $Y = X \smallsetminus (C \cup \lbrace a\rbrace)$ sont liés
à $a$ et chacun à exactement deux sommets de $C$, formant un côté ; deux points
de $Y$ sont liés si et seulement s'ils correspondent à des côtés opposés, et
les deux paires de côtés opposés de $C$ correspondent aux deux triangles de
sommet $a$ (pages 81 et 82)\footnote{La page écrit « l'ombre sur $C$ d'un pt de
$Y$ est un côté de $C$ ». La phrase de la page 81, « est lié exactement à un
point de $C$ diag $\lbrace x, x'\rbrace$ », se poursuit page 82 par l'autre
diagonale : un point de chaque diagonale, donc deux sommets formant un côté.}. La page 82
écrit en travers un diagramme d'ensembles d'hexagones qu'on ne restitue pas.

Page 83, une question : quels graphes finis ont une « enveloppe » pavée ?
Exemples dessinés : le point, le segment, le carré, l'hexagone, le pentagone,
deux segments, le sablier. Pages 84 et 85 : un \emph{anti-carré} (somme de deux
segments) ; $C \mapsto X \smallsetminus (C \cup \lbrace \text{centre}\rbrace)$ est une
bijection des carrés sur les anti-carrés ; et le pavé engendré par un carré
vérifie une propriété universelle évidente pour les plongements dans un complexe
de la famille,
\[
\mathrm{Pl}(\mathrm{Pav}(C), \mathcal{X}) \xrightarrow{\ \sim\ }
\mathrm{Pl}(C, \mathcal{X}),
\]
de sorte que la donnée d'un carré de $\mathcal{X}$ équivaut à celle d'un pavé
pointé ; $\mathrm{Pav}(C)^\circ \simeq \mathrm{Pav}(C^*)$, $C^*$ le carré
dual\footnote{Une partie notable des pages 84 et 85 est illisible ; on ne
garde que les énoncés écrits en formules.}.

\pagerange{86}{90}
\section*{18. Produits, coproduits et graphe opposé (pages 86 à 90)}

Le \emph{produit} de graphes $(\Gamma_i)$ a pour sommets $\prod \Gamma_i$, deux
sommets étant liés s'ils diffèrent en une seule coordonnée, où ils sont liés
(c'est le produit cartésien des graphes, et non leur produit catégorique). Le
graphe \emph{opposé} $\Gamma^\circ$ a les mêmes sommets, deux sommets distincts
étant liés si et seulement s'ils ne le sont pas dans $\Gamma$ (le graphe
complémentaire) ; $\Gamma$ est \emph{auto-opposé} s'il est isomorphe à
$\Gamma^\circ$. Le \emph{coproduit} est $\overset{\circ}{\coprod} \Gamma_i =
\bigl(\prod \Gamma_i^\circ\bigr)^\circ$.

\textbf{Proposition} (page 86). Un pavé est auto-opposé ; de façon équivalente,
le coproduit de deux triades est un pavé. C'est vrai : le pavé est le produit
de deux triangles, son opposé est le coproduit de deux triades (les triades
étant les opposés des triangles), et le pavé est isomorphe à son
complémentaire\footnote{La page l'énonce avec un point d'interrogation. Le
pavé est le graphe de Paley d'ordre 9, ou le graphe de la grille $3 \times 3$,
qui est auto-complémentaire ; le dossier 69 (pages 73 à 75) établit le même
fait. Noms de nous.}.

Pages 87 et 88 : le coproduit d'un doublet et d'un graphe discret est un graphe
bi-discret ; le coproduit de deux doublets est la somme de deux segments (le
carré dessiné réduit à ses diagonales), et le produit de deux segments est le
carré\footnote{La page dessine le carré sous « Coproduit des deux segments ».
Avec la définition de la page 86, le coproduit de deux segments est le graphe
complet à quatre sommets ; c'est leur \emph{produit} qui est le carré. La
suite de la page (« le graphe complété semble \ldots ») paraît s'en apercevoir.}.
En mettant sur l'un des facteurs une structure de doublet et sur l'autre une
structure de segment, on obtient au total les trois structures de carré sur
l'ensemble à quatre éléments $\lbrace a, b\rbrace \times \lbrace a', b'\rbrace$,
qui correspondent à ses trois partitions de type $(2, 2)$. Un carré est donc la
donnée d'un ensemble à quatre éléments et d'une telle partition $(C_\alpha)_{\alpha
\in D}$ (ses diagonales) ; une orientation est donnée par des bijections
$\varphi_\alpha : C_\alpha \to C_\beta$ (« passer au sommet suivant ») avec
$\varphi_\beta \varphi_\alpha = -\mathrm{id}$, $-$ désignant l'involution qui
échange les deux sommets d'une diagonale ; d'où $\mathrm{or}(C) \simeq D^* \times
D$\footnote{Lecture incertaine du symbole devant $C_\alpha$ et de la légende
« carré dual » sous l'accolade.}.

Les pages 89 et 90, au crayon et en largeur, dressent le grand diagramme des
ensembles de carrés et d'hexagones d'un pavé, avec les groupes ($\mathfrak{S}_3$,
$\mathfrak{S}_2$, $D_3$, $D_6$, $\mathbb{F}_2 \times \mathbb{F}_2$) ; une marge
y note « invariants cube et hypercube dans polyèdres et polytopes ; prismes
triangulaires, pentagonaux, hexagonaux », programme que reprendront les pages
134 et 135. On ne restitue pas ces diagrammes.

\pagerange{92}{100}
\section*{19. Complexes cubiques généralisés (pages 92 à 100)}

La liasse qui commence page 92 porte un titre : « Complexes cubiques
généralisés (complexes \emph{triadiques} ou \emph{trilatères} ?) ». C'est le
tournant du dossier.

Soit $S$ un graphe (sans boucles ni arêtes multiples), $A \subset
\mathfrak{P}_2(S)$ ses arêtes, $T \subset \mathfrak{P}_3(S)$ l'ensemble de ses
triangles\footnote{La page écrit $t \in \mathfrak{P}_3(A)$ ; il faut
$\mathfrak{P}_3(S)$, comme la transcription le note.}.

\textbf{Axiome 1.} Toute arête est contenue dans un unique triangle.

\textbf{Axiome 2.} Pour tout triangle $t$ et tout sommet $s \notin t$, $s$ est lié à
un sommet de $t$ (unique d'après l'axiome 1).

Leur conjonction est l'\emph{axiome (C)}. Ce sont exactement les axiomes d'un
\emph{quadrangle généralisé} dont les droites ont trois points : les sommets
sont les points, les triangles les droites ; « un point hors d'une droite est
collinéaire à un unique point de cette droite »\footnote{Nom de nous. Les
quadrangles généralisés ont été introduits par Tits (1959) ; ceux d'ordre $(2,
t)$ n'existent, à part les cas dégénérés, que pour $t = 1, 2, 4$, et chacun est
unique : la grille $3 \times 3$, le quadrangle $W(2)$ à 15 points, et le
quadrangle à 27 points des droites de la surface cubique. Le résultat est
classique (on le trouve dans la monographie de Payne et Thas, \emph{Finite
generalized quadrangles}, 1984) ; cité de mémoire. Ici $t = c$.}.

\textbf{Autour d'un triangle} (pages 92 à 97). Supposons l'axiome 1, $A \neq
\emptyset$, et fixons un triangle $t_0 = \lbrace a, b, c\rbrace$ pour lequel
l'axiome 2 est vrai. Pour $\alpha \in t_0$, soit $\widetilde E_\alpha$
l'ensemble des sommets hors de $t_0$ liés à $\alpha$. Alors
\[
(1) \qquad S \simeq t_0 \amalg \coprod_{\alpha \in t_0} \widetilde E_\alpha .
\]
Les arêtes sont \emph{basiques} (dans $t_0$), \emph{verticales} (de $\alpha$ à
$\widetilde E_\alpha$) ou \emph{horizontales supérieures} (dans $S \smallsetminus
t_0$). L'axiome 1 donne successivement :

\begin{itemize}
\item \textbf{Prop.~1.} Pour $s \in \widetilde E_\alpha$, il y a un unique $s' \in
  \widetilde E_\alpha$ lié à $s$ (le troisième sommet du triangle de l'arête
  $\lbrace \alpha, s\rbrace$). D'où une involution sans point fixe $\sigma_\alpha$
  de $\widetilde E_\alpha$, et l'involution $\sigma = \sigma_{t_0}$ de $S$ égale à
  l'identité sur $t_0$ et à $\sigma_\alpha$ sur $\widetilde E_\alpha$\footnote{La
  page numérote d'abord cet énoncé « 2 » (lecture incertaine) puis y renvoie
  comme « Prop.~1 » ; elle écrit $\alpha \to \alpha'$ pour $s \mapsto s'$.}.
\item \textbf{Prop.~2.} Si $s, s' = \sigma_\alpha s$ et $t \in \widetilde
  E_\beta$, $\beta \neq \alpha$, alors $t$ n'est pas lié à la fois à $s$ et à
  $s'$.
\item \textbf{Prop.~3.} Un triangle contenant une arête « transverse » (entre
  $\widetilde E_\alpha$ et $\widetilde E_\beta$, $\alpha \neq \beta$) a un sommet
  dans chacun des trois $\widetilde E_\alpha$. Ces triangles \emph{supérieurs}
  forment un ensemble $\widetilde T$, avec $\tilde\varphi : \widetilde T \to
  \prod_\alpha \widetilde E_\alpha$.
\item \textbf{Prop.~4.} Pour toute arête basique $u = \lbrace \alpha,
  \beta\rbrace$, la projection $\widetilde T \to \widetilde E_\alpha \times
  \widetilde E_\beta$ est injective, d'image l'ensemble des couples liés.
\end{itemize}

\textbf{Scholie} (page 96). Moyennant l'axiome 1 et l'axiome 2 pour $t_0$, le
graphe se résume à la donnée de a) un ensemble $t_0$ à trois éléments, b) une
famille $(\widetilde E_\alpha)_{\alpha \in t_0}$ d'ensembles munis d'involutions
sans point fixe $\sigma_\alpha$, c) une partie $\widetilde T \subset \widetilde P
= \prod_\alpha \widetilde E_\alpha$, soumise aux conditions (*) les projections
sur les produits de deux facteurs sont injectives, (**) si $(s, t)$ est dans
l'image, $(\sigma_\alpha s, t)$ n'y est pas\footnote{Une troisième condition,
(***), écrite en travers de la marge et tournée de 90°, n'est lue qu'en
partie ; c'est l'axiome que la page 101 appellera (Ka).} ; les liaisons sont
celles qu'on devine (2°) $\alpha$)--$\delta$)), et les triangles sont
\[
T \simeq \lbrace t_0\rbrace \amalg \coprod_{\alpha \in t_0} E_\alpha \amalg
\widetilde T, \qquad E_\alpha = \widetilde E_\alpha / \sigma_\alpha
\]
(le triangle de base, les triangles verticaux $\lbrace \alpha\rbrace \cup u$,
$u$ une orbite de $\sigma_\alpha$, et les triangles supérieurs). « Rien
n'empêche pour l'instant qu'on ait $\widetilde T = \emptyset$ ! »

\textbf{L'axiome 2 pour les autres triangles} (pages 98 à 100). Pour un
triangle vertical $\lbrace \alpha, s, s'\rbrace$ il donne la \textbf{Prop.~5} :
tout $t \in \widetilde E_\beta$ ($\beta \neq \alpha$) est lié à exactement l'un
de $s$, $s'$. D'où trois corollaires : si $s, t$ sont liés, $s'$ et $t' =
\sigma_\beta t$ le sont ; $\widetilde T$ est stable par l'involution $\prod
\sigma_\alpha$ de $\widetilde P$ ; l'ensemble $\widetilde E_\beta(s)$ des
voisins de $s$ dans $\widetilde E_\beta$ est une section de $\widetilde E_\beta
\to E_\beta$, et « être liés » est une bijection $\widetilde E_\beta(s) \to
\widetilde E_\gamma(s)$. Par suite, si $\widetilde E_\alpha \neq \emptyset$,
$\widetilde E_\beta$ et $\widetilde E_\gamma$ ont même cardinal, et si de plus
celui-ci n'est pas nul les trois $\widetilde E_\alpha$ ont même cardinal $2c$.
Les cas où certains $\widetilde E_\alpha$ sont vides donnent des graphes
dégénérés (un triangle seul, une « étoile » de triangles en un sommet), que
l'\textbf{axiome 3} exclut : les $\widetilde E_\alpha$ sont non vides. Sous
l'axiome 3, l'axiome 2 pour les triangles supérieurs est alors automatique : la
\textbf{Prop.~6}, barrée sur la page mais lisible, le démontre (si $s_1 \in
\widetilde E_a$ n'était lié à aucun des sommets $t, u$ d'un triangle supérieur
$\lbrace s, t, u\rbrace$, $\sigma_a s_1$ serait lié aux deux, contredisant la
Prop.~2)\footnote{L'énoncé de l'axiome 3 sur la page (« on n'est pas dans le
cas \ldots\ où un sommet est dans au moins 2 triangles ») est d'une lecture
incertaine ; on donne la forme « les $\widetilde E_\alpha$ sont non vides », qui
le précède. La « Résolution » qui suit, « la donnée d'un complexe satisfaisant
les axiomes (1), (2) \ldots\ équivaut à celle de », s'interrompt au bas de la
page 100.}.

\pagerange{101}{105}
\section*{20. La réduction à deux problèmes (pages 101 à 105)}

La page 101 reformule la Scholie : un ensemble $t_0$ à trois éléments, un
ensemble non vide $\widetilde\Gamma$ (les triangles supérieurs) avec une
involution sans point fixe $\sigma_{\widetilde\Gamma}$, et trois quotients
$\widetilde\Gamma \to \widetilde E_\alpha$ compatibles aux involutions, avec

\begin{enumerate}
\item[$\alpha$)] les involutions induites $\sigma_\alpha$ sont sans point fixe ;
\item[$\beta_{\widetilde\Gamma}$)] pour $\alpha \neq \beta$, $\widetilde\Gamma
  \to \widetilde E_\alpha \times \widetilde E_\beta$ est injective, et pour toute
  orbite $u$ de $\sigma_\alpha$ et tout $t \in \widetilde E_\beta$, un et un seul
  $s \in u$ a $(s, t)$ dans l'image ;
\item[$\gamma$)] l'axiome (Ka), énoncé plus haut sur des familles $(s_\alpha)$,
  $(t_\alpha)$ et qu'on ne peut restituer\footnote{L'axiome (Ka) est la condition
  (***) de la marge de la page 96, lue en partie seulement.}.
\end{enumerate}

Passons aux quotients : $\Gamma = \widetilde\Gamma/\sigma$, $E_\alpha =
\widetilde E_\alpha/\sigma_\alpha$. Alors $\widetilde\Gamma$ s'identifie, pour
chaque $\alpha$, à l'image inverse du revêtement double $\widetilde E_\alpha \to
E_\alpha$ par $\Gamma \to E_\alpha$, et la donnée équivaut à : un ensemble
$\Gamma$ et trois quotients $\varphi_\alpha : \Gamma \to E_\alpha$, trois
revêtements doubles $\widetilde E_\alpha \to E_\alpha$, et un système transitif
d'isomorphismes entre leurs trois images inverses sur $\Gamma$. L'axiome
$\beta_{\widetilde\Gamma}$ équivaut alors à

\textbf{$\beta_\Gamma$)} pour $\lbrace \alpha, \beta\rbrace \in \mathfrak{P}_2(t_0)$,
$\Gamma \to E_\alpha \times E_\beta$ est \emph{bijective}.

C'est mot pour mot l'axiome I de la page 1 : $\Gamma \subset \prod_\alpha
E_\alpha$ est un carré latin d'ordre $c$, et $\operatorname{card}\Gamma = c^2$,
$\operatorname{card}\widetilde\Gamma = 2c^2$\footnote{Le rapprochement avec la
page 1 est de nous ; les pages 101 à 103 ne renvoient pas en arrière.}. La page
102 l'écrit : « la classification des cas "complexes cubiques généralisés" se
décompose en deux questions exactement » :

\begin{enumerate}
\item[I)] décrire les systèmes de trois quotients $\Gamma \to \Gamma_\alpha$
  d'un ensemble non vide vérifiant $\beta_\Gamma$ --- les carrés latins ;
\item[II)] pour un tel système, classer ses « revêtements d'ordre 2 », puis
  ceux qui vérifient l'axiome (Ka).
\end{enumerate}

Le problème I) est celui des pages 1 à 16 ; le problème II), dans le cas où le
carré latin est la table d'un groupe commutatif $V$ (c'est-à-dire $\Gamma =
\operatorname{Ker}(V^3 \to V)$), est celui des extensions centrales des pages 32
à 39. La page 116 montrera que $V$ est nécessairement un
$\mathbb{F}_2$-espace vectoriel, et la page 39 donne l'extension pour $\dim V =
2$\footnote{Ce double rapprochement est notre lecture ; c'est aussi celle que
suggère le diagramme de la page 121 (section 23). Le dossier 69 fait le même
chemin depuis le graphe cubique : lemme 7, puis ses pages 14 et 15, puis
107 à 113.}.

\textbf{Les nombres} (page 104). Avec $c = \operatorname{card} E_a$ :
\[
\begin{aligned}
&\operatorname{card} S = 3(1 + 2c), && \operatorname{card} E_S(s) = 2(1 + c), \\
&\operatorname{card} A = 3(1+c)(1+2c), && \operatorname{card} T = (1+c)(1+2c),
\end{aligned}
\]
chaque sommet étant sur $1 + c$ triangles, et $\operatorname{card}\widetilde
\Gamma = 2c^2$, $\operatorname{card}\Gamma = c^2$, $\operatorname{card}\widetilde
E_a = 2c$. Tout est juste. La page 105 classe les cas dégénérés (points isolés,
un triangle, une étoile de triangles en un sommet), et s'arrête sur « Si un
sommet porte au moins 2 triangles »\footnote{Plusieurs mots entourés y sont
illisibles.}.

\pagerange{106}{116}
\section*{21. Automorphismes des graphes trialitaires (pages 106 à 116)}

Le titre est de sa main : « Automorphismes des graphes trialitaires » (le mot
« complexes » y est biffé et remplacé par « graphes »). Soit $(S, A)$ un graphe
vérifiant les axiomes 1 à 3, $c \geq 1$.

Pour un triangle $\underline t = \lbrace a, b, c\rbrace$, la \emph{symétrie}
$\sigma_{\underline t}$ est la bijection de $S$ égale à l'identité sur
$\underline t$ et qui envoie $x \in \widetilde E_a$ sur l'unique élément de
$\widetilde E_a$ lié à $x$ (c'est l'involution $\sigma_{t_0}$ de la section
19).

\textbf{Proposition} (page 106). a) $\sigma_{\underline t}$ est un automorphisme
de $(S, A)$, dont l'ensemble des points fixes est $\underline t$. b)
L'application $\underline t \mapsto \sigma_{\underline t}$ est injective, et son
image est formée des automorphismes involutifs $u$ dont l'ensemble des points
fixes est un triangle et qui stabilisent chaque triangle passant par un de ces
points fixes\footnote{La page énonce b) avec la seule condition « $u$ involutif et
$\operatorname{card} S^u = 3$ », et ajoute « C'est évident ». Ainsi énoncé, b) est
faux. Pour $c = 1$, la transposition du pavé $3 \times 3$ est une involution
dont les points fixes sont les trois points de la diagonale, qui forment une
triade et non un triangle. Pour $c = 2$, dans le modèle des paires de $\lbrace
1, \ldots, 6\rbrace$, la double transposition $(34)(56)$ fixe exactement les trois
paires $12$, $34$, $56$, qui forment un triangle, sans être la symétrie de ce
triangle (qui est la triple transposition $(12)(34)(56)$) : il y a 45 telles
involutions pour 15 triangles. La condition supplémentaire qu'on a ajoutée
suffit : un tel $u$ préserve chaque paire $\lbrace s, s'\rbrace$ d'un triangle
$\lbrace \alpha, s, s'\rbrace$, et, sans point fixe hors de $\underline t$, il
l'échange ; il coïncide donc avec $\sigma_{\underline t}$.}. Pour $g \in
\mathcal{G} = \operatorname{Aut}(S, A)$ on a $g \sigma_{\underline t} g^{-1} =
\sigma_{g(\underline t)}$, de sorte que le groupe $\mathcal{G}^\circ$ engendré par
les symétries est distingué dans $\mathcal{G}$.

\textbf{Théorème} (page 107). $\mathcal{G}^\circ$ est transitif sur l'ensemble des
arêtes orientées (couples de sommets liés), donc sur les sommets et sur les
triangles. \textbf{Corollaire} : les $\sigma_{\underline t}$ forment une seule
classe de conjugaison de $\mathcal{G}^\circ$. La marge excepte le cas $c = 1$ :
pour le pavé, $\mathcal{G}^\circ \simeq \mathfrak{S}_3 \times \mathfrak{S}_3$
est d'indice 2 dans $\mathcal{G} \simeq \mathfrak{S}_2 \cdot (\mathfrak{S}_3
\times \mathfrak{S}_3)$, transitif sur les sommets mais non sur les arêtes (les
symétries des lignes et celles des colonnes ne sont pas conjuguées dans
$\mathcal{G}^\circ$)\footnote{La fin de l'énoncé du théorème est perdue dans des
insertions. On a vérifié par machine la transitivité sur les arêtes orientées
pour $c = 2$ et $c = 4$, et les ordres de $\mathcal{G}^\circ$ : 36, 720 et
25\,920 pour $c = 1, 2, 4$. Pour $c = 4$, $\mathcal{G}^\circ$ est donc
d'indice 2 dans $\mathcal{G} = W(E_6)$ ; les pages ne le disent pas, et ce que
dit la page 115 (« $\mathcal{G} = \mathcal{G}^\circ$ ») ne concerne que $c =
2$.}.

La démonstration repose sur deux lemmes. \textbf{Lemme 1} : a) si $s, t$ sont
liés, $t = \sigma_{\underline t} s$ pour un triangle $\underline t$ (n'importe
quel autre triangle passant par le troisième sommet $u$ du triangle $\lbrace s,
t, u\rbrace$ ; il en existe dès que $c \geq 1$) ; b) si $s, t$ ne sont pas liés,
ils ont un voisin commun, et $t$ est image de $s$ par un produit de deux
symétries. \textbf{Lemme 2} : si $s, t$ sont deux voisins distincts et non liés d'un
même sommet $a$, il existe un produit de deux symétries qui fixe $a$ et échange
$s$ et $t$ ; et (page 108) si $c \geq 2$ on peut le choisir induisant
l'identité sur un triangle $\underline t_0 = \lbrace a, b, c\rbrace$ ne contenant ni
$s$ ni $t$. La « recette » des pages 108 et 109 construit ce produit : avec $u$ un
voisin commun de $s$ et $t$ hors de $a$, lié à $b$, les triangles $\lbrace s, u,
v\rbrace$ et $\lbrace t, u, w\rbrace$, et les triangles $\underline\sigma = \lbrace c, v,
v'\rbrace$ et $\underline\tau = \lbrace c, w, w'\rbrace$, on a $\sigma_{\underline\sigma}
s = u$, $\sigma_{\underline\tau} u = t$, les deux symétries échangent $a$ et $b$
et fixent $c$, et
\[
W = \sigma_{\underline\tau} \sigma_{\underline\sigma}
\]
envoie $(a, s)$ sur $(a, t)$ en induisant l'identité sur $\lbrace a, b,
c\rbrace$\footnote{La lettre $\underline\sigma$ nommant un triangle est d'une
lecture incertaine. Une première version de la recette (page 108) et un NB
(page 109, continué page 112) sont barrés.}.

\textbf{Transitivité sur les triangles supérieurs} (pages 110 à 114). Le
corollaire de la page 110 (barré dans sa moitié inférieure) affirme que les
produits pairs de symétries de triangles verticaux, qui fixent $a$, $b$, $c$,
opèrent transitivement sur $\widetilde\Gamma$. Le \textbf{Lemme 1} de la page 111,
valable pour $c$ fini $\geq 3$ (« c'est faux si $c = 1, 2$ »), affirme que deux
éléments $s, t$ de $\widetilde E_a$ se joignent par un $s_1 \in \widetilde E_a$ et
des $u, v \in \widetilde E_b$ avec $u$ lié à $s$ et $s_1$, $v$ à $s_1$ et $t$ ;
la démonstration, par l'absurde et par cardinalité, sépare $\widetilde E_a$ en
deux parties $A'$, $A''$ et aboutit (page 113) à $\widetilde E_c(r_i) =
\complement\widetilde E_c(r_i)$, absurde. Le \textbf{Lemme 2} et son corollaire
(pages 113 et 114) introduisent, pour un sommet $b$ de $\underline t_0$, le groupe
$\mathcal{G}_b$ engendré par les produits $\sigma_{v_1}\sigma_{v_2}$ de deux
symétries de triangles verticaux en $b$ : c'est un $\mathbb{F}_2$-espace vectoriel
qui respecte $a$, $b$, $c$, et il est transitif sur $\widetilde E_a$ (par des
produits de quatre symétries) ; un « corollaire plus fort », barré, l'affirmait
transitif sur $\widetilde\Gamma$\footnote{Les indices de $\pi =
\sigma_{w_1}\sigma_{v_2}$ dans le lemme 2 sont surchargés ; lecture incertaine.}.

\textbf{Ordres} (page 115). Le nombre de triangles épinglés (triangle muni d'un
ordre total sur ses sommets) est le nombre d'arêtes orientées, $6(1+c)(1+2c)$ ;
celui des bitriangles épinglés est $2c^2$ fois plus grand, $12c^2(1+c)(1+2c)$.
Pour $c = 2$, un automorphisme qui induit l'identité sur un bitriangle épinglé
est l'identité, et $\mathcal{G} = \mathcal{G}^\circ$ est simplement transitif sur
les bitriangles épinglés : il est d'ordre $12 \cdot 4 \cdot 5 \cdot 3 = 720$. Pour
$c = 1$, $\mathcal{G}$ est d'ordre $12 \cdot 1 \cdot 3 \cdot 2 = 72$. En
général le nombre de bitriangles est $\frac{1}{72}$ du nombre des bitriangles
épinglés, $\frac16 c^2(1+c)(1+2c)$. Tout cela est juste.

\textbf{Le cas commutatif} (page 116). Pour la relation $xyz = 1$ sur un groupe
$G$ et un élément $a$, l'application
\[
\varphi_a : (x, y, z) \longmapsto (x,\ a^{-1}z^{-1},\ a^{-1}y^{-1})
\]
préserve la relation pour tous $(x, y, z)$ seulement si $a^2 = 1$ (prendre $y = z
= 1$), puis $(z^{-1}a)^2 = 1$ pour tout $z$ : tout élément est d'ordre 2, $G$ est
commutatif, c'est un espace vectoriel sur $\mathbb{F}_2$, et la condition est
alors satisfaite. Écrit additivement, $\varphi_c : (x, y, z) \mapsto (x, z + c, y +
c)$ et $\varphi_a \varphi_b : (x, y, z) \mapsto (x, y + c, z + c)$ avec $c = a +
b$\footnote{Ici $c$ est un élément du groupe, non le cardinal $c$ ; la page
réutilise la lettre.}. Ce sont les symétries des triangles verticaux, vues sur le
carré latin $\Gamma$ : le sous-groupe distingué $\mathcal{G}_a$ engendré par les
produits pairs en $a$ opère trivialement sur $\widetilde E_a$, transitivement sur
$\widetilde E_b$ et $\widetilde E_c$, et il est isomorphe au groupe des parties
paires de $E_a$ ; $\mathcal{G}$ opère transitivement sur $\widetilde\Gamma \simeq
\mathcal{G}/\mathfrak{g}$. Voilà pourquoi le carré latin du problème I) doit être la
table d'un $\mathbb{F}_2$-espace vectoriel, et $c$ une puissance de 2.

\pagerange{117}{123}
\section*{22. La tour et ses tableaux (pages 117 à 123)}

Soit $\mathcal{T}_\beta$, $\beta \in \lbrace -1, 0, 1, 2\rbrace$, les graphes de la
famille, avec $c_\beta = 0, 1, 2, 4$. Pour $\beta \leq \gamma$ on note
$c_{\beta\gamma}$ le nombre de sous-structures $\mathcal{T}_\beta$ dans un
$\mathcal{T}_\gamma$, $c'_{\beta\gamma}$ le nombre de sous-structures
\emph{épinglées}, et $N_\beta = c'_{\beta\beta}$ l'ordre du groupe des
automorphismes de $\mathcal{T}_\beta$ ; on a
\[
c'_{\beta\gamma} = N_\beta\, c_{\beta\gamma}, \qquad c'_{\beta\gamma} =
c'_{\alpha\gamma}\, \nu_{\alpha\beta\gamma},
\]
$\nu_{\alpha\beta\gamma}$ étant le nombre de prolongements d'une structure
épinglée $\mathcal{T}_\alpha$ en une structure épinglée $\mathcal{T}_\beta$.
On trouve (page 118)
\[
\nu_{-1,0,\gamma} = 2c_\gamma^2, \qquad \nu_{0,1,\gamma} = c_\gamma - 1, \qquad
\nu_{1,2,\gamma} = c_\gamma - 2,
\]
d'où, pour $\gamma$ convenable,
\[
\begin{aligned}
c'_{-1,\gamma} &= 6(1+c_\gamma)(1+2c_\gamma), &
c'_{0,\gamma} &= 12c_\gamma^2(1+c_\gamma)(1+2c_\gamma), \\
c'_{1,\gamma} &= 12c_\gamma^2(c_\gamma^2-1)(2c_\gamma+1), &
c'_{2,\gamma} &= 12c_\gamma^2(c_\gamma^2-1)(c_\gamma-2)(2c_\gamma+1),
\end{aligned}
\]
et les deux tableaux, qui sont justes :
\[
\begin{array}{c|cccc}
c'_{\beta\gamma} & \gamma = -1 & 0 & 1 & 2 \\ \hline
\beta = -1 & 6 & 36 & 90 & 270 \\
0 & & 72 & 720 & 8640 \\
1 & & & 720 & 25\,920 \\
2 & & & & 51\,840
\end{array}
\qquad
\begin{array}{c|cccc}
c_{\beta\gamma} & \gamma = -1 & 0 & 1 & 2 \\ \hline
\beta = -1 & 1 & 6 & 15 & 45 \\
0 & & 1 & 10 & 120 \\
1 & & & 1 & 36 \\
2 & & & & 1
\end{array}
\]
avec
\[
c_{0,\gamma} = \tfrac{1}{72}\,c'_{0,\gamma} = \tfrac16\, c_\gamma^2 (c_\gamma+1)
(2c_\gamma+1), \qquad c_{1,\gamma} = \tfrac{1}{720}\, c'_{1,\gamma} =
\tfrac{1}{60}\, c_\gamma^2 (c_\gamma^2 - 1)(2c_\gamma+1) .
\]\footnote{La page 119
écrit $c_{0,\gamma} = \frac{1}{72} c_{0,\gamma}$ (sans prime) et y ajoute un
facteur $(c_\gamma - 1)$, qui donnerait 360 au lieu de 120 pour $c = 4$ ; le
tableau suit la bonne formule. Les formules fermées de la page 120,
$\frac16 c(c+1)(2c+1)$ et $\frac{1}{30}c(c+1)(2c+1)(c-1)$, perdent chacune un
facteur $c$ : la première donnerait 5 au lieu de 10 pour $c = 2$, la seconde 18
au lieu de 36 pour $c = 4$ (la transcription le note pour la seconde).}
Les ordres $N_\beta$ sont $6$, $72$, $720$, $51\,840 = 2^7 \cdot 3^4 \cdot 5$,
que la page 122 recalcule en travers d'une copie d'examen ($128 \times 405 =
51\,840$).

\textbf{Le tableau de la page 120.} Il récapitule, pour $c = 0, 1, 2, 4$ :
triangles par sommet $c + 1 = 1, 2, 3, 5$ ; voisins $2(c+1) = 2, 4, 6, 10$ ;
sommets $3 + 6c = 3, 9, 15, 27$ ; arêtes $3, 18, 45, 135$ ; triangles $1, 6, 15,
45$, dont $3c$ verticaux et $2c^2 = 0, 2, 8, 32$ transversaux ; sommets non liés à
un sommet donné $4c = 0, 4, 8, 16$, qu'il écrit $2^{\nu+2}$ ; et l'ordre du groupe
d'automorphismes, $6$, $72$, $720$, $51\,840$\footnote{La case $c = 2$ de l'ordre
est surchargée (un « 360 » recouvert). La dernière ligne du dernier tableau de la
page (72, 360, 2160) ne s'accorde pas avec les $c'_{0,\gamma}$ de la page 118
(72, 720, 8640) ; on suit la page 118, que confirment les comptes directs.}.
Il note enfin « $c = 2^\nu$ » : $2c^2 = 2^{2\nu+1}$ pour $\nu = -\infty, 0, 1,
2$.

\textbf{Catégories de paires} (page 123). La catégorie des $\mathcal{T}_\alpha$
munis d'une sous-structure $\mathcal{T}_\beta$ est équivalente à :
\[
\begin{array}{c|cccc}
\beta \backslash \alpha & -1 & 0 & 1 & 2 \\ \hline
\text{rien} & (\mathrm{Ens})_3 & (\mathrm{Ens})_{3,3} & (\mathrm{Ens})_6 &
\ast \\
-1 & (\mathrm{Ens})_3 & (\mathrm{Ens})_3 \times (\mathrm{Ens})_2 &
(\mathrm{Ens})_{2,2,2} & \text{« Trialité »} \\
0 & & (\mathrm{Ens})_{3,3} & (\mathrm{Ens})_{3,3} & (\mathrm{Ens})_3 \times
(\mathrm{Ens})_{3,3} \\
1 & & & (\mathrm{Ens})_6 & (\mathrm{Ens})_6 \times (\mathrm{Ens})_2
\end{array}
\]
On a contrôlé chaque case sur les ordres : le groupe d'automorphismes de la
catégorie doit être d'ordre $N_\alpha / c_{\beta\alpha}$, et c'est le
cas\footnote{La case $(0, 2)$ porte sur la page $(\mathrm{Ens})_{3,3} \times
(\mathrm{Ens})_{3,3}$, d'ordre $72^2$ ; il faut $51\,840/120 = 432 = 6 \cdot 72$,
soit $(\mathrm{Ens})_3 \times (\mathrm{Ens})_{3,3}$, ce que donne la page 52
(section 11). La case $\ast$ ($\mathcal{T}_2$ seul) est laissée sans
équivalent : $W(E_6)$ n'est pas le groupe d'une structure aussi simple.}. Le mot
« Trialité » pour un graphe cubique muni d'un triangle est remarquable : le
stabilisateur d'un triangle est d'ordre $51\,840/45 = 1152$, l'ordre du groupe de
Weyl de $F_4$, extension par $\mathfrak{S}_3$ (les permutations des trois sommets
du triangle) de $W(D_4)$, et c'est la trialité de $D_4$ qui opère\footnote{Le
dossier 69 (section 9) établit l'isomorphisme avec $W(F_4)$. Le rapprochement
avec la trialité est de nous ; la page n'écrit que le mot.}. Le second tableau de la
page, pour les chaînes $\mathcal{T}_\alpha \subset \mathcal{T}_\beta \subset
\mathcal{T}_\gamma$, est lu sur un aperçu insuffisant et toutes ses entrées sont
douteuses\footnote{Avec les comptes ci-dessus, les groupes des chaînes $(-1, 0,
1)$, $(-1, 0, 2)$, $(-1, 1, 2)$, $(0, 1, 2)$ et $(-1, 0, 1, 2)$ sont d'ordres 12,
72, 96, 144 et 24. Les entrées lues $(\mathrm{Ens})_{2,2,2} \times
(\mathrm{Ens})_2$ (96) et $(\mathrm{Ens})_{3,3} \times (\mathrm{Ens})_2$ (144)
s'accordent ; les trois autres lectures ($(\mathrm{Ens})_3$, $(\mathrm{Ens})_3
\times (\mathrm{Ens})_3$, $(\mathrm{Ens})_3 \times (\mathrm{Ens})_2$) non.}.

\pagerange{121}{125}
\section*{23. Revêtements, torseurs et relations entre symétries (pages 121, 124, 125)}

La page 121, de brouillon, dessine le diagramme qui réunit les deux moitiés du
dossier : $V \xrightarrow{\varphi_i} \Gamma \xrightarrow{p_i} E_i/2 = V$, le relèvement
$\psi_i : V \to \widetilde\Gamma$, l'extension $\lbrace \pm 1\rbrace \to
\widetilde\Gamma \to \Gamma$, et le revêtement $\widetilde\Gamma \to \widetilde
E_i$\footnote{Lecture du diagramme incertaine : la flèche $p_i$ et le but $E_i/2$
sont surchargés, $\psi_i$ est en pointillé.}. C'est la situation des pages 32 à
39 appliquée au problème II) des pages 101 à 103. On y lit aussi les relations
$\tau_2\tau_1 = \tau_3\tau_2 = \tau_1\tau_3 = \mathrm{id}$ de la page 10, et des
calculs de commutateurs de symétries, $\sigma_x \sigma_y \sigma_x^{-1}\sigma_y^{-1}
= \sigma_{\sigma_x y}\sigma_y^{-1}$, conséquence de $g\sigma_{\underline t}g^{-1}
= \sigma_{g \underline t}$.

Page 124 : pour trois triangles $t_0$, $t$, $t'$ convenablement placés, les
symétries agissent sur six sommets $x, y, z, x', y', z'$ par
\[
\sigma_{t_0} : x \leftrightarrow x', \ y \leftrightarrow y', \ z \leftrightarrow
z' ; \quad \sigma_t : x \leftrightarrow y', \ y \leftrightarrow z', \ z
\leftrightarrow x' ; \quad \sigma_{t'} : x \leftrightarrow z', \ y
\leftrightarrow x', \ z \leftrightarrow y' ,
\]
et $\sigma_{t_0}\sigma_t = \sigma_t\sigma_{t'} = \sigma_{t'}\sigma_{t_0} = \pi$ est
d'ordre 3 : les trois symétries engendrent un $\mathfrak{S}_3$. Page 125 : avec
un triangle supérieur $\lbrace s, t, u\rbrace$, l'identité encadrée
\[
\sigma_t \sigma_s \sigma_u = \sigma_{t_0} \quad \text{sur } \widetilde E_a
\smallsetminus \lbrace s, s'\rbrace ,
\]
où $\sigma_s$ désigne la symétrie du triangle vertical passant par $s$, et ses
réécritures, dont $\sigma_t \sigma_s \sigma_t^{-1} = \sigma_{\sigma_t s}$ ; la marge
dit « short-cut !! »\footnote{La page 124 évoque aussi l'action de
$\mathfrak{S}_I$ et de $\mathfrak{S}_I \times \mathfrak{S}_J$ sur les trois blocs
$\xi_1\xi_2\xi_3 / \xi^*_4\xi^*_5\xi^*_6$, etc. ; le passage est en grande partie
illisible. La page 125 note aussi $\frac12(3+6c)(2+2c) = 3(1+2c)(1+c)$, le nombre
d'arêtes.}.

\pagerange{127}{129}
\section*{24. Un sommet marqué et l'espace $\mathbb{F}_2^I$ (pages 127 à 129)}

Soit $(S, A)$ un graphe $T_\nu$, $\nu \in \lbrace 0, 1, 2\rbrace$, $c = 2^\nu$, et $a
\in S$. Soit $I$ l'ensemble des $c + 1$ triangles de sommet $a$, $X =
\coprod_{i \in I} X_i$ l'ensemble des voisins de $a$ ($X_i$ les deux autres
sommets du triangle $i$, torseur sous $\mathbb{F}_2$), et $E = \prod_{i \in I}
X_i$, torseur sous $V = \mathbb{F}_2^I$. Soit $Y$ l'ensemble des $4c$ sommets non
liés à $a$. Tout $x \in Y$ est lié à un unique élément $\varphi_i(x)$ de chaque
$X_i$, d'où $\varphi : Y \to E$.

Les symétries $\sigma_i$ des triangles $i \in I$ opèrent sur $E$ : $\sigma_{i_0}$
échange les deux points de chaque $X_i$, $i \neq i_0$, et fixe $X_{i_0}$. Si
$\eta : V \to \mathbb{F}_2$ est la somme des coordonnées, $\delta$ la diagonale et
$z = \delta(1)$, on a donc $(\sigma_g)_E = \rho(g + \delta\eta(g))$, où $\rho$ est
l'action naturelle de $V$ sur $E$. Soit $u(x) = x + \delta\eta(x)$. Son noyau est
$\lbrace 0\rbrace$ si $\operatorname{card} I$ est pair et $\lbrace 0, z\rbrace$ si
$\operatorname{card} I$ est impair ; or $\operatorname{card} I = 1 + 2^\nu$ est
pair pour $\nu = 0$ et impair pour $\nu = 1, 2$. Donc :

\begin{itemize}
\item pour $\nu = 0$ ($c = 1$, le pavé), $u$ est un isomorphisme ; $\varphi$
  commute aux actions, elle est injective, et comme $\operatorname{card} Y = 4 =
  2^{c+1} = \operatorname{card} E$, elle est bijective ;
\item pour $\nu = 1, 2$, $u$ a pour noyau $\lbrace 0, z\rbrace$ et pour image le
  sous-espace $V' = \operatorname{Ker}\eta$ des vecteurs pairs, de dimension
  $c$\footnote{La page conclut au contraire « si $\nu \neq 0$, $u$ est
  injectif », après avoir bien établi que $z$ est dans le noyau exactement quand
  $\operatorname{card} I$ est impair ; c'est un lapsus, et la suite de la page
  (« $\operatorname{card} E = 2^{c+1} = 2^2 = 4$, donc $\varphi$ est bijectif »)
  ne vaut que pour $c = 1$. Elle écrit en outre « cas du point » pour $\nu = 0$.
  Le b) qui suit, pour $\nu = 1, 2$, s'interrompt sur « $\operatorname{Ker}\eta =
  V$ \ldots », qu'on lit $V'$.}.
\end{itemize}

Les pages s'arrêtent là pour $\nu \geq 1$. Ce qu'on peut dire, et qu'on a
vérifié : pour $c = 4$, $\varphi$ est injective et son image est une orbite de
$V'$, de sorte que les seize non-voisins de $a$ forment un torseur sous $V'$, de
dimension 4 --- c'est ce qu'affirme la page 135 ; pour $c = 2$, en revanche,
$\varphi$ est deux-à-un sur quatre points (dans le modèle des paires de
$\lbrace 1, \ldots, 6\rbrace$, avec $a = 12$, les non-voisins $13$ et $23$ ont les mêmes
projections), et la description ne s'étend pas telle quelle\footnote{Pour $c
= 4$, dans le modèle de Schläfli avec $a = a_1$, un non-voisin est codé par la
partie de $\lbrace 2, \ldots, 6\rbrace$ des plans tritangents $\lbrace a_1, b_j,
c_{1j}\rbrace$ où il rencontre $c_{1j}$ plutôt que $b_j$ : $a_k \mapsto \lbrace
k\rbrace$, $c_{jk} \mapsto \lbrace 2, \ldots, 6\rbrace \smallsetminus \lbrace j,
k\rbrace$, $b_1 \mapsto \lbrace 2, \ldots, 6\rbrace$ ; ce sont exactement les seize
parties de cardinal impair. Le dossier 69 (pages 16 et 17) code de la même
façon les non-voisins par des sections.}.

\pagerange{130}{131}
\section*{25. Les mêmes axiomes en termes de triangles (pages 130 et 131)}

Titre de sa main : « Graphes pseudotrialitaires en termes de triangles ». On se
donne un ensemble $X$ de « sommets » et $T \subset \mathfrak{P}_3(X)$, et l'on
définit une arête comme une paire contenue dans un triangle. Axiomes : a) deux
triangles distincts ont au plus un point commun (toute arête est dans un unique
triangle) ; b) trois sommets deux à deux joints par des arêtes forment un triangle
de $T$ ; c) pour $\underline t \in T$ et $a \notin \underline t$, il existe un
unique triangle passant par $a$ et rencontrant $\underline t$. C'est l'axiome (C)
de la page 92 dans le langage des points et des droites. La condition d)
« chaque sommet est sur trois triangles » fait de $X$ une \emph{configuration}
(chaque droite a trois points et chaque point est sur trois droites) ; avec
$X'$ l'ensemble des droites et $R$ l'incidence, $X' \to \mathfrak{P}_3(X)$ et $X
\to \mathfrak{P}_3(X')$ sont injectives, et la page entreprend de vérifier que le
dual $X'$ satisfait à son tour a) à d) : d) est immédiat, a) est clair, b) se
démontre par l'absurde, c) est énoncé\footnote{La lecture « conf. » est
incertaine. Avec d), on est dans le cas $c = 2$ : la configuration est le
quadrangle $W(2)$ (configuration de Cremona--Richmond $15_3$), qui est
effectivement auto-dual. Nom de nous.}.

\pagerange{132}{135}
\section*{26. Diagrammes $A_n$, prismes et carrés (pages 132 à 135)}

Les dernières pages décrivent des configurations du graphe cubique, chacune
par la catégorie des données équivalentes, dans le style des pages 17 à 39 du
dossier 69 ; on vérifie chaque fois la cohérence avec $|W(E_6)| = 51\,840$.

\begin{itemize}
\item \textbf{Page 132.} Un cube dont les sommets portent des $\lambda_{ij}$ ;
  « cube $\leftrightarrow$ cube ponctué » ; $(\mathrm{Ens})_{2,2,2}$ et
  $(\mathrm{Ens})_2 \times (\mathrm{Ens})_4$ ; $\mathfrak{S}_3 \cdot (\pm 1)^3
  \simeq \mathfrak{S}_2 \times \mathfrak{S}_4$, isomorphisme exact (groupe
  hyperoctaédral d'ordre 48).
\item \textbf{Diagramme $A_2$} (page 133), qui équivaut à un « triangle
  ponctué » (une arête, et son triangle) : la catégorie des graphes cubiques
  munis d'un tel diagramme est équivalente à celle des ensembles à huit éléments
  munis d'une involution sans point fixe, ou des 4-cubes, ou des torseurs sous
  $\mathfrak{S}_4 \cdot (\pm 1)^4 \simeq W(B_4)$, le groupe de Weyl de $SO(9)$,
  d'ordre 384 ; et $51\,840/384 = 135$, le nombre d'arêtes.
\item \textbf{Diagramme $A_3$} : $(\mathrm{Ens})_2 \times (\mathrm{Ens})_4$,
  d'ordre 48, soit $1080$ diagrammes (les chemins induits $x - y - z$, $x$ et $z$
  non liés : $27 \cdot 10 \cdot 8 / 2 = 1080$).
\item \textbf{Diagramme $A_4$} : « pavé $+$ triade et triangle sécants »,
  « hexagone $+$ clou centré », $(\mathrm{Ens})_2 \times (\mathrm{Ens})_3$, d'ordre
  12, soit 4320 ; \textbf{diagramme $A_5$} : « hexagone ponctué », « triade
  épinglée », groupe d'ordre 12 ($720 \cdot 6 = 4320$ triades
  épinglées)\footnote{L'indice de « $\mathrm{Ens}_{2,4}$ » est surchargé, un 4
  repris en 2 ; avec $(\mathrm{Ens})_2 \times (\mathrm{Ens})_3$ l'ordre est le
  bon. La page 66 annonçait déjà « $A_4 \Leftrightarrow A_5 \Leftrightarrow P_6
  \Leftrightarrow$ bitriangle épinglé ». Le sens précis de « diagramme $A_n$ »
  n'est pas défini sur ces feuillets ; les nombres sont ceux des configurations
  que le dossier 69 appelle balises et préprismes.}.
\item \textbf{Prismes triangulaires} (page 134) : « pavé $+$ triangle »,
  $(\mathrm{Ens})_2 \times (\mathrm{Ens})_3 \times (\mathrm{Ens})_3$, d'ordre 72,
  soit 720 prismes (le nombre du dossier 69). \textbf{Bicarrés} : $(\mathrm{Ens})_2
  \times (\mathrm{Ens})_2 \times (\mathrm{Ens})_2 \simeq \mathrm{Tors}((\pm 1)^3)$,
  et $\frac12 (720 \times 6 \times 3) = 2^4 \cdot 3^4 \cdot 5 = 51\,840/8$.
\item \textbf{Page 135.} Une bande de trois carrés n'existe pas dans un graphe
  cubique (le troisième sommet de deux triangles serait forcé d'être égal à un
  autre) ; d'où : pas de prismes à base pentagonale ou hexagonale, mais des
  prismes à base carrée, c'est-à-dire des cubes. Enfin, la catégorie des graphes
  cubiques à sommet marqué $a$ équivaut à celle des couples formés d'un ensemble
  $I$ à cinq éléments (les triangles en $a$) et d'un torseur $T$ sous $V_I =
  (\mathbb{F}_2^I)'$, l'espace de dimension 4 engendré par des $e_i$ de seule
  relation $\sum e_i = 0$ ; $T$ est l'ensemble des sommets non liés à $a$
  (groupe d'ordre $120 \cdot 16 = 1920 = |W(D_5)|$, et $51\,840/1920 = 27$). Les
  carrés de centre $a$ correspondent aux parties $I' \subset I$ de cardinal 2 et
  aux points de $T/V_{(I')}$, et la catégorie des graphes cubiques munis d'un
  carré équivaut à $(\mathrm{Ens})_{2,2} \times (\mathrm{Ens})_3$\footnote{Le
  crochet ouvert à cet endroit n'est pas refermé : l'argument continue au-delà du
  dossier. L'ordre $8 \cdot 6 = 48$ donne $1080$ carrés, le nombre de 4-cycles
  induits que compte le dossier 69.}.
\end{itemize}

C'est sur cette description autour d'un sommet --- un ensemble à cinq éléments et
un torseur sous un espace de dimension 4 sur $\mathbb{F}_2$ --- que le dossier
s'arrête : elle rejoint le graphe de la page 18, construit sur un
$\mathbb{F}_2$-espace vectoriel, par lequel les notes sur les vingt-sept avaient
commencé.

\end{document}
